% Transcribed from CMM110 Final (Lectures 13-21).docx.
\IfFileExists{Book\space NL/template/kaobook.cls}{
  \IfFileExists{main.tex}{\documentclass[main.tex]{subfiles}}{
    \IfFileExists{../../main.tex}{\documentclass[../../main.tex]{subfiles}}{
      \IfFileExists{../main.tex}{\documentclass[../main.tex]{subfiles}}{
        \documentclass[Book\space NL/main.tex]{subfiles}
      }
    }
  }
}{
  \IfFileExists{../../template/kaobook.cls}{\documentclass[../../main.tex]{subfiles}}{\documentclass{book}}
}
\IfFileExists{Book\space NL/template/kaobook.cls}{}{\IfFileExists{../../template/kaobook.cls}{}{\InputIfFileExists{review-notes-standalone.tex}{}{\InputIfFileExists{Book\space NL/appendices/appendix 4-Review Notes/review-notes-standalone.tex}{}{}}}}
\begin{document}
\ifSubfilesClassLoaded{
  \ifcsname appendixchapterparent\endcsname
  \else
    \standaloneappendixchaptersetup{D}{4}
  \fi
}{}
\IfFileExists{\booknlroot appendices/appendix 4-Review Notes/images/d4-figure-glut4-present-secretory-vesicles-cytoplasm.png}{
  \edef\reviewnotesimagepath{\booknlroot appendices/appendix 4-Review Notes/images/}
}{
  \edef\reviewnotesimagepath{\booknlroot appendices/appendix 4-Review Notes/images/}
}
\chapter{Metabolism, Molecular Biology, and Genetics Review}
\labch{appendix-d4-metabolism-molecular-biology-genetics-review}

\setcounter{margintocdepth}{1}
\margintoc
\section{Additional Resources}

Metabolic pathway relational maps and databases (searchable):

https://mededucation.stanford.edu/wp-content/uploads/2024/01/FullSubwayMap\_V1023\_Web.pdf

https://www.kegg.jp/pathway/map01100

https://www.kegg.jp/kegg/pathway.html\#metabolism

https://smpdb.ca/view

\section{Carbohydrate Metabolism}

\subsection{The Human Metabolism}

Raw materials are needed for the functioning of the human body. These materials are often referred to as 'nutrients.’ These nutrients are processed by the metabolism, a never-ending series of biochemical reactions that degrade and process food molecules into energy and the building blocks for biosynthesis (through catabolic processes), or synthesize various molecules for use within cells (through anabolic processes) to repair and build new cells.

All cells in the human body perform metabolic processes, and require energy and raw materials for use in those processes. Metabolic processes may either utilize or generate energy. Breakdown (catabolic) processes often result in energy generation, while synthesis (anabolic) processes require the expenditure of energy in order to synthesize end products. The most major form of cellular energy is ATP (Adenosine Triphosphate), generated by the process of cellular respiration.

Metabolic pathways often make use of enzymes in order to perform reactions (without those enzymes, the activation energy of the reaction would result in no spontaneous and rapid reaction). Enzymes are able to accelerate reactions by decreasing the activation energy required to perform reactions. Often, in the metabolic process, Many enzymes are utilized in long chains of processes to break down large molecules in the metabolic pathway efficiently.

Key enzyme classifications are listed in Table~\ref{tab:appendix-table-d4-enzyme-reaction-classes}.

\pagelayout{wide}
\begin{center}
\captionof{table}{Functional classes of enzymes and representative reactions.}\labtab{appendix-table-d4-enzyme-reaction-classes}
\small
\renewcommand{\arraystretch}{1.15}
\begin{tabularx}{\textwidth}{p{0.18\textwidth}p{0.28\textwidth}X}
\toprule
\textbf{Enzyme class} & \textbf{General reaction} & \textbf{Representative subclasses and details} \\
\midrule
Oxidoreductases & Oxidation--reduction reactions. & \textbf{Dehydrogenases}: addition or removal of hydrogen.\newline \textbf{Oxidases}: transfer two electrons to oxygen, forming hydrogen peroxide, or to one-half oxygen, forming water.\newline \textbf{Oxygenases}: incorporate oxygen into the product.\newline \textbf{Hydroxylases}: incorporate one-half oxygen into the product as --OH and form water.\newline \textbf{Peroxidases}: use hydrogen peroxide as an oxygen donor, forming water. \\

Transferases & Transfer a chemical group from one substrate to another. & \textbf{Kinases}: transfer phosphate from ATP to a substrate. \\

Hydrolases & Hydrolyse C--O, C--N, O--P, and C--S bonds. & Examples include esterases, proteases, phosphatases, and deamidases. \\

Lyases & Add groups across a carbon--carbon double bond. & Examples include dehydratases, hydratases, and decarboxylases. \\

Isomerases & Catalyse intramolecular rearrangements. & \\

Ligases (synthetases) & Form bonds between two substrates. & Often coupled to ATP utilisation, with an intermediate phosphorylated enzyme or substrate. \\
\bottomrule
\end{tabularx}
\end{center}
\pagelayout{margin}

\subsection{Micro and Macronutrients}

Carbohydrates are organic compounds composed mainly of carbon, hydrogen, and oxygen. Many simple carbohydrates have an approximate 1:2:1 ratio of carbon, hydrogen, and oxygen, although this is not universal. Carbohydrates provide energy and carbon skeletons for biosynthesis. Excess energy intake can contribute to weight gain, while very low carbohydrate intake may lead to ketosis; metabolic acidosis is not an inevitable result of carbohydrate deficiency.

Dietary carbohydrates are digested to monosaccharides before absorption. Common disaccharides include sucrose (glucose and fructose), maltose (two glucose units), and lactose (glucose and galactose). Absorbed monosaccharides can be oxidized for energy, stored as glycogen, or converted to other molecules, including fatty acids, depending on metabolic state.

\subsection{Transport of Glucose}

Because glucose is hydrophilic, facilitated diffusion through membrane glucose transporters is required for cellular uptake. GLUT1 and GLUT3 support basal uptake in many tissues, with GLUT3 especially important in neurons. GLUT2 is a high-capacity transporter expressed in liver and pancreatic beta cells, among other sites. GLUT4 is insulin-responsive in skeletal muscle and adipose tissue and is also expressed in cardiac muscle. Other GLUT family members have distinct tissue distributions and substrates. Table~\ref{tab:appendix-table-d4-glucose-transporter-tissues} summarizes the transporter data shown in the source.

\pagelayout{wide}
\begin{center}
\captionof{table}{Glucose transporters in humans.}\labtab{appendix-table-d4-glucose-transporter-tissues}
\small
\renewcommand{\arraystretch}{1.15}
\begin{tabularx}{\textwidth}{p{0.11\textwidth}Xp{0.12\textwidth}p{0.31\textwidth}}
\toprule
\textbf{Transporter} & \textbf{Tissues expressed} & \textbf{$K_t$ (mM)\textsuperscript{*}} & \textbf{Reported role\textsuperscript{†}} \\
\midrule
GLUT1 & Ubiquitous & 3 & Basal glucose uptake \\

GLUT2 & Liver, pancreatic islets, intestine & 17 & In liver and kidney, removal of excess glucose from blood; in pancreas, regulation of insulin release \\

GLUT3 & Brain (neuronal), testis (sperm) & 1.4 & Basal glucose uptake \\

GLUT4 & Muscle, fat, heart & 5 & Activity increased by insulin \\

GLUT5 & Intestine (primarily), testis, kidney & 6\textsuperscript{‡} & Primarily fructose transport \\

GLUT6 & Spleen, leukocytes, brain & $>5$ & Possibly no transport function \\

GLUT7 & Small intestine, colon & 0.3 & --- \\

GLUT8 & Testis & $\sim 2$ & --- \\

GLUT9 & Liver, kidney & 0.6 & --- \\

GLUT10 & Heart, lung, brain, liver, muscle, pancreas, kidney & 0.3\textsuperscript{§} & --- \\

GLUT11 & Heart, skeletal muscle, kidney & 0.16 & --- \\

GLUT12 & Skeletal muscle, heart, prostate, small intestine & --- & --- \\
\bottomrule
\end{tabularx}
\end{center}
\noindent\textbf{[Source note: The table image includes footnote markers for $K_t$, role, GLUT5, and GLUT10, but their definitions are not visible in the supplied image.]}
\pagelayout{margin}

In the absence of insulin stimulation, much of the GLUT4 in muscle and adipose cells is stored in intracellular vesicles. Insulin receptor signaling promotes translocation and fusion of these vesicles with the plasma membrane, increasing glucose transport into the cells. Muscle contraction can also stimulate GLUT4 translocation in skeletal muscle through insulin-independent signaling.

\begin{center}\includegraphics[width=0.82\textwidth,height=0.68\textheight,keepaspectratio]{\reviewnotesimagepath d4-figure-glut4-present-secretory-vesicles-cytoplasm.png}\par\captionof{figure}{Insulin signaling promotes GLUT4 translocation to the plasma membrane, increasing glucose uptake in insulin-responsive tissues.}\labfig{appendix-figure-d4-glut4-present-secretory-vesicles-within-cytoplasm}\end{center}

Facilitative glucose transporters alternate between outward-facing and inward-facing conformations. Glucose binding on one side of the membrane shifts the transporter to expose the binding site to the other side, where glucose is released; the unoccupied transporter then returns to its initial conformation.

\begin{center}\includegraphics[width=0.82\textwidth,height=0.68\textheight,keepaspectratio]{\reviewnotesimagepath d4-figure-glucose-transporters-conformations-conformation-binding.png}\par\captionof{figure}{Alternating outward- and inward-facing conformations mediate facilitated glucose transport.}\labfig{appendix-figure-d4-glucose-transporters-conformations-conformation-binding-site}\end{center}

\subsection{Fates of Glucose}

Glucose supports energy production through glycolysis and supplies ribose-5-phosphate and NADPH through the pentose phosphate pathway. It also provides carbon for glycogen synthesis and for biosynthetic pathways, including glycoprotein and extracellular-matrix carbohydrate synthesis. Humans store glucose primarily as glycogen; starch and sucrose are dietary carbohydrates, not human storage forms.

\begin{center}\includegraphics[width=0.82\textwidth,height=0.68\textheight,keepaspectratio]{\reviewnotesimagepath d4-figure-glucose-utilized-energy-generation-inside.png}\par\captionof{figure}{Overview of glucose use in energy production, biosynthesis, and storage.}\labfig{appendix-figure-d4-glucose-utilized-energy-generation-inside-glycolysis}\end{center}

Glucose is central to several metabolic pathways. During fasting, hepatic glycogenolysis helps maintain blood glucose, while gluconeogenesis synthesizes glucose from lactate, glycerol, and glucogenic amino acids. Glycogenolysis releases glucose-1-phosphate from glycogen; the liver can convert glucose-6-phosphate to free glucose, whereas muscle uses it locally. When glucose is abundant, glycogenesis stores it as glycogen, mainly in liver and skeletal muscle.

\subsection{Entry of Common Substances Into Glycolysis}

Glycolysis metabolizes glucose and other monosaccharides after they are converted to pathway intermediates. Dietary starch and disaccharides must first be digested to absorbable monosaccharides; intracellular glycogen can enter glycolysis after glycogenolysis.

Dietary starch is hydrolyzed by amylases to shorter oligosaccharides and maltose, then by intestinal brush-border enzymes to monosaccharides, including glucose, which can enter glycolysis.

During fasting or exercise, glycogen phosphorylase cleaves alpha-1,4 linkages in endogenous glycogen by phosphorolysis, producing glucose-1-phosphate, which phosphoglucomutase converts to glucose-6-phosphate. A debranching enzyme handles alpha-1,6 branch points and releases a small amount of free glucose.

Other hexoses enter glycolysis after conversion to phosphorylated intermediates. Sucrase hydrolyzes sucrose to glucose and fructose. In the liver, fructokinase converts fructose to fructose-1-phosphate, which aldolase B splits into glyceraldehyde and dihydroxyacetone phosphate; triokinase and triose-phosphate isomerase convert these products to glycolytic intermediates. In other tissues, hexokinase can convert fructose to fructose-6-phosphate.

Mannose may be phosphorylated by hexokinase, resulting in mannose-6-phosphate. M6P is then isomerized by phosphomannose isomerase, resulting in fructose-6-phosphate.

The common pathways are showcased in the diagram on the following page.

\begin{center}\includegraphics[width=0.82\textwidth,height=0.68\textheight,keepaspectratio]{\reviewnotesimagepath d4-figure-glycolysis-common-pathways.png}\par\captionof{figure}{Entry points from common dietary carbohydrates into glycolysis.}\labfig{appendix-figure-d4-glycolysis-common-pathways}\end{center}

\subsection{Glycolysis}

\paragraph{Overall reaction and products}

\begin{center}\includegraphics[width=0.82\textwidth,height=0.68\textheight,keepaspectratio]{\reviewnotesimagepath d4-figure-glucose-breakdown-oxidation-products.png}\par\captionof{figure}{Overview of glucose oxidation through glycolysis and its principal products.}\labfig{appendix-figure-d4-glucose-breakdown-oxidation-products}\end{center}



Glycolysis (glyco meaning glucose, and lysis meaning breakdown of) is the breakdown of organic 'fuels’ (mainly glucose, but some other carbohydrates may be broken down in this manner, for instance, glycogen and starch, which may be broken down into phosphorylated glucose, or other mono or disaccharides capable of being broken down into appropriate metabolic intermediates) which occurs in the cytosol of cells. It is a 10 step process which has two 5 step phases: the preparatory (investment) phase and the payoff phase. In the preparatory phase, glucose is phosphorylated and subsequently converted into glyceraldehyde-3-phosphate (G3P) through a breakdown process. This phase utilizes ATP, and results in the precursor for the payoff phase. In the payoff (oxidative) phase, G3P is converted into pyruvate, and the end products (ATP and NADH) of glycolysis are formed. Overall, 6-carbon glucose is broken down into 3-carbon pyruvate, and a net gain of 2 ATP and 2 NADH occurs. This net reaction is as follows: Glucose + 2 NAD\textsuperscript{+} + 2ADP + 2P\textsubscript{i} \(\rightarrow\) 2 pyruvate + 2 NADH + 2H\textsuperscript{+} + 2 ATP + 2H\textsubscript{2}O. Note that all carbon in the end product is present in the pyruvate, which enters into the citric acid cycle.

In the preparatory (energy investment) phase, glucose is first fixed by phosphorylation using the \textbf{hexokinase} \textbf{(kinases transfer phosphate groups)} enzyme. ATP is used in this stage (one phosphate group is transferred to the 6th carbon in glucose), as such this reaction is called the \textbf{first priming reaction}. The product of this reaction is \textbf{glucose-6-phosphate}, which is then converted by \textbf{phosphohexose isomerase} into fructose-6-phosphate by isomerization (from 6-carbon central ring to 5-carbon central ring and a side chain). Fructose-6-phosphate then undergoes a \textbf{second priming reaction}, where \textbf{phosphofructokinase} adds another phosphate group to convert fructose-6-phosphate into \textbf{fructose 1,6-bisphosphate}. After the ATP-utilizing priming reactions, fructose 1,6 bisphosphate is split into two molecules and isomerized by \textbf{aldolase }(resulting in \textbf{glyceraldehyde 3 phosphate} and \textbf{dihydroxyacetone phosphate}). Each product of this reaction is a 3 carbon compound. Dihydroxyacetone phosphate is then isomerized by triose phosphate isomerase into glyceraldehyde 3 phosphate (a process called 'interconversion'). This results in a total conversion of each fructose 1,6 bisphosphate into two molecules of glyceraldehyde 3 phosphate (which is often referred to as the \textbf{common product}). \textbf{In summary, }two ATP molecules are utilized as energy in the priming phase of the glycolysis process. In this process, glucose is converted into 2 molecules of the 3 carbon glyceraldehyde 3 phosphate through isomerization and breakdown processes. Note that the phosphorylations are the \textbf{rate-limiting} step of glycolysis, meaning that they are the slowest steps in the process.

In the \textbf{payoff phase}, the energy carrier ATP and the electron carrier NADH (from NAD\textsuperscript{+}) are synthesized from their precursors (in reality, a more appropriate but technically incorrect form is 'filled up'), fueled by the breakdown of glyceraldehyde-3-phosphate into pyruvate (the end product of glycolysis). In the payoff phase, a cytosolic (inorganic) phosphate group is attached to each of the two glyceraldehyde-3-phosphate molecules, resulting in the oxidation of a total (each G-3-P produces 1 NADH) of 2 NAD\textsuperscript{+} into 2 NADH molecules (with 2H\textsuperscript{+} as an additional product). This is done by the enzyme \textbf{glyceraldehyde-3-phosphate dehydrogenase}, and the reaction is referred to as dehydrogenation. The end product of this phosphorylation is\textbf{ 1,3-bisphosphoglycerate }(a phosphate is added to the first carbon).

Following this, two (total) ATP are produced from ADP during the \textbf{substrate level phosphorylation} process, which is assisted by \textbf{phosphoglycerate kinase}. In this process, a phosphate is removed from 1,3-diphosphoglycerate and added to ADP to form 1 ATP per molecule. This results in the formation of the molecule \textbf{3-phosphoglycerate}. This reaction is also known as the 'first ATP forming reaction,’ since it results in the first ATP production of the glycolysis pathway.

Following the formation of ATP, 3-phosphoglycerate is converted by phosphoglycerate mutase into \textbf{2-phosphoglycerate} when the phosphate bond is moved to the second carbon. 2-phosphoglycerate is then rearranged into \textbf{phosphoenolpyruvate} When \textbf{enolase} removes one water molecule from each of the two 2-phosphoglycerate molecules in a lysis reaction. Finally, phosphoenolpyruvate then undergoes a second substrate-level phosphorylation, where \textbf{pyruvate kinase }removes the remaining phosphate and produces one ATP per molecule of phosphoenolpyruvate. As such, this results in the production of \textbf{pyruvate}. This marks the end of the glycolysis pathway. Note that pyruvate undergoes further breakdown in the citric acid cycle, and the various glycolytic intermediates may be utilized in various other metabolic processes.

The phases of glycolysis are summarized in the diagram on the following page (key phases and enzymes are highlighted).


\subsection{Oxidation of Pyruvate}

However, metabolic processes do not end at the formation of pyruvate, as it may be further broken down in the process of pyruvate oxidation. This breakdown process occurs in the mitochondria (which contains the requisite enzymes), necessitating the transport of pyruvate into the mitochondrial matrix as glycolysis occurs in the cytosol. This transport is done by transport proteins. Following transport into the mitochondria, a multienzyme (3 enzymes with cofactors, referred to in general as the pyruvate dehydrogenase complex (PDC)) complex catalyzes three reactions (\textbf{pyruvate decarboxylation}). In the first reaction, pyruvate's carboxyl group is removed and released as CO\textsubscript{2}. Following this, the remaining two-carbon segment Is oxidized, forming acetate, which releases electrons to NAD\textsuperscript{ +}, forming NADH. Coenzyme A is then attached, forming acetyl CoA. This process is showcased in the diagram below.

\begin{center}\includegraphics[width=0.82\textwidth,height=0.68\textheight,keepaspectratio]{\reviewnotesimagepath d4-figure-however-metabolic-end-formation-pyruvate.png}\par\captionof{figure}{Major metabolic fates of pyruvate after glycolysis.}\labfig{appendix-figure-d4-however-metabolic-end-formation-pyruvate-further}\end{center}

The following diagrams showcase the role of various enzymes and essential vitamins (cofactors) in the process of pyruvate decarboxylation involving the removal of carbon dioxide, transfer of an electron to NAD\textsuperscript{+}, and the addition of acetyl CoA. Coenzyme forms derived from water-soluble vitamins are listed in Table~\ref{tab:appendix-table-d4-water-soluble-vitamin-coenzymes}.

\begin{center}\includegraphics[width=0.82\textwidth,height=0.68\textheight,keepaspectratio]{\reviewnotesimagepath d4-figure-enzymes-essential-vitamins-cofactors-pyruvate.png}\par\captionof{figure}{Enzymes and vitamin-derived cofactors involved in pyruvate oxidation.}\labfig{appendix-figure-d4-enzymes-essential-vitamins-cofactors-pyruvate-decarboxylation}\end{center}

\begin{center}\includegraphics[width=0.82\textwidth,height=0.68\textheight,keepaspectratio]{\reviewnotesimagepath d4-figure-enzymes-essential-vitamins-cofactors-pyruvate-02.png}\par\captionof{figure}{Enzymes and vitamin-derived cofactors involved in pyruvate oxidation.}\labfig{appendix-figure-d4-enzymes-essential-vitamins-cofactors-pyruvate-decarboxylation-02}\end{center}

\begin{center}\includegraphics[width=0.82\textwidth,height=0.68\textheight,keepaspectratio]{\reviewnotesimagepath d4-figure-enzymes-essential-vitamins-cofactors-pyruvate-03.png}\par\captionof{figure}{Enzymes and vitamin-derived cofactors involved in pyruvate oxidation.}\labfig{appendix-figure-d4-enzymes-essential-vitamins-cofactors-pyruvate-decarboxylation-03}\end{center}

\subsection{Alternate Fates of Pyruvate}

When oxygen is present as the last electron acceptor, complete oxidation occurs. However, under hypoxic and anaerobic conditions, pyruvate may be transformed into lactate by lactate fermentation. Typically, this occurs most In skeletal muscles during vigorous contractions (AKA exercise) and in hypoxic tumor cells (yay warburg effect again!). Lactate fermentation occurs because NADH cannot be reoxidized into NAD\textsuperscript{+}, which means that it is unable to receive electrons from glycolysis and pyruvate oxidation. As such, pyruvate ferments, resulting in the regeneration of NAD\textsuperscript{+}. This is often referred to as lactic acid fermentation, which utilizes the enzyme lactate dehydrogenase. The oxidation of NADH is due to giving electrons to lactate (one electron per molecule of lactate). This process is shown below.

\begin{center}
\captionof{table}{Coenzyme derivatives of water-soluble vitamins.}\labtab{appendix-table-d4-water-soluble-vitamin-coenzymes}
\small
\renewcommand{\arraystretch}{1.15}
\begin{tabularx}{\textwidth}{p{0.40\textwidth}X}
\toprule
\textbf{Vitamin} & \textbf{Coenzyme form} \\
\midrule
Vitamin B\textsubscript{1} (thiamine) & Thiamine pyrophosphate (TPP) \\

Vitamin B\textsubscript{2} (riboflavin) & Flavin mononucleotide (FMN); flavin adenine dinucleotide (FAD) \\

Vitamin B\textsubscript{3} (pantothenic acid) & Coenzyme A (CoA) \\

Vitamin B\textsubscript{5} (niacin) & Nicotinamide adenine dinucleotide (NAD); nicotinamide adenine dinucleotide phosphate (NADP) \\

Vitamin B\textsubscript{6} (pyridoxine) & Pyridoxal phosphate (PLP); pyridoxamine phosphate (PAMP) \\

Vitamin B\textsubscript{7} (biotin) & Biocytin \\

Vitamin B\textsubscript{9} (folic acid) & Tetrahydrofolic acid (THFA) \\

Vitamin B\textsubscript{12} (cyanocobalamin) & Deoxyadenosyl cobalamin \\

Vitamin C (ascorbic acid) & Not known \\
\bottomrule
\end{tabularx}
\end{center}
% CONTENT AUDIT: Preserve the vitamin-number assignments as transcribed; verify the source diagram because pantothenic acid and niacin appear to have swapped B-vitamin numbers.
\noindent\textbf{[Verify against a reliable reference: the source labels pantothenic acid as B\textsubscript{3} and niacin as B\textsubscript{5}; those labels appear transposed.]}

\begin{center}\includegraphics[width=0.82\textwidth,height=0.68\textheight,keepaspectratio]{\reviewnotesimagepath d4-figure-oxygen-present-last-electron-acceptor-02.png}\par\captionof{figure}{Aerobic oxidation links the citric acid cycle to the respiratory chain.}\labfig{appendix-figure-d4-oxygen-present-last-electron-acceptor-complete-02}\end{center}

\subsection{The Citric Acid/Krebs Cycle}

The citric acid (TCA) cycle has two purposes: the oxidation of Acetyl-CoA into CO\textsubscript{2} to produce energy, and the creation of precursors utilized in various biosynthesis reactions (including carbohydrate, lipid, amino acid, nucleotide and porphyrin molecules).

The metabolic intermediates of the TCA cycle are linked to many metabolic processes, such as serving as a way to process (preprocessed) degraded amino acids, or \textbf{gluconeogenesis} (the production of glucose). Furthermore, these intermediates may also be utilized in the production or breakdown of fatty acids, proteins, or hemes (through porphyrin). key metabolic intermediates and their link to the citric acid cycle are shown below. This results in the TCA often being referred to as an amphibolic (containing both catabolic and anabolic processes) cycle.

\begin{center}\includegraphics[width=0.82\textwidth,height=0.68\textheight,keepaspectratio]{\reviewnotesimagepath d4-figure-metabolic-intermediates-tca-cycle-linked.png}\par\captionof{figure}{Citric acid-cycle intermediates connect carbohydrate, lipid, and amino-acid metabolism.}\labfig{appendix-figure-d4-metabolic-intermediates-tca-cycle-linked-many}\end{center}

During each iteration of the TCA cycle, one \textbf{Acetyl CoA} is imported, one ATP is generated through substrate level phosphorylation, and a number of electrons are gained (and later utilized to generate ATP through the electron transport chain). The electrons gained are carried by NADH and FADH\textsubscript{2} to the \textbf{electron transport chain}.

\textbf{(focus on energy-generating cycle components!)}

One of the (somewhat directly) ATP-generating reactions is the substrate level phosphorylation of \textbf{succinyl COA} into \textbf{succinate} (through the conversion of GTP produced by the reaction to GDP). After this reaction, succinate releases an electron to FAD\textsuperscript{+}, resulting in FADH\textsubscript{2} and \textbf{fumarate}, which is subsequently rearranged by the addition of a water molecule into \textbf{malate}.

The \textbf{net equation} for the citric acid cycle is: acetyl CoA + 3 NAD\textsuperscript{+} + FAD\textsuperscript{+} + ADP + HPO\textsubscript{4}\textsuperscript{2-} \(\rightarrow\)2 CO\textsubscript{2} + CoA + 3 NADH + FADH\textsubscript{2} + ATP (total ATP from the cycle + pyruvate oxidation is 12, 9 from 3 NADH, 2 from one FADH\textsubscript{2 }and one from direct generation, as seen in the diagram below). An overview of this process is shown on the following page.

\begin{center}\includegraphics[width=0.82\textwidth,height=0.68\textheight,keepaspectratio]{\reviewnotesimagepath d4-figure-net-equation-citric-acid-cycle.png}\par\captionof{figure}{Net reaction of one turn of the citric acid cycle per acetyl-CoA.}\labfig{appendix-figure-d4-net-equation-citric-acid-cycle-acetyl}\end{center}

\begin{center}\includegraphics[width=0.82\textwidth,height=0.68\textheight,keepaspectratio]{\reviewnotesimagepath d4-figure-net-equation-citric-acid-cycle-02.png}\par\captionof{figure}{Net reaction of one turn of the citric acid cycle per acetyl-CoA.}\labfig{appendix-figure-d4-net-equation-citric-acid-cycle-acetyl-02}\end{center}

\subsection{Gluconeogenesis}

During fasting, gluconeogenesis helps maintain blood glucose by synthesizing glucose from lactate, glycerol, and glucogenic amino acids. It occurs primarily in the liver and, during prolonged fasting, the renal cortex. Fatty-acid oxidation supplies energy for gluconeogenesis, but even-chain fatty acids do not provide net glucose in humans.

Gluconeogenesis shares several reversible reactions with glycolysis but bypasses its three strongly regulated, effectively irreversible steps. Fructose-1,6-bisphosphatase converts fructose-1,6-bisphosphate to fructose-6-phosphate; glucose-6-phosphatase converts glucose-6-phosphate to glucose. Pyruvate is converted to oxaloacetate by pyruvate carboxylase, then to phosphoenolpyruvate by phosphoenolpyruvate carboxykinase (PEPCK).

\begin{center}\includegraphics[width=0.82\textwidth,height=0.68\textheight,keepaspectratio]{\reviewnotesimagepath d4-figure-gluconeogenesis-quite-similar-glycolysis-however.png}\par\captionof{figure}{Gluconeogenesis bypasses the irreversible reactions of glycolysis using distinct enzymes.}\labfig{appendix-figure-d4-gluconeogenesis-quite-similar-glycolysis-however-three}\end{center}

In more detail, the first bypass reaction (at the 10th step of glycolysis) entails the conversion of pyruvate into phosphoenolpyruvate through two exergonic (ATP consuming) reactions. This reaction may occur either in the cytosol or mitochondria, depending on the availability of precursors. For pyruvate, the reaction occurs inside of the mitochondria. This entails either biosynthesis from alanine transamination or importation of pyruvate (the two most predominant precursors along with lactate, which utilizes a separate pathway).

Pyruvate carboxylase is a mitochondrial, biotin-dependent enzyme that uses ATP to convert pyruvate to oxaloacetate. PEPCK then uses GTP to convert oxaloacetate to phosphoenolpyruvate; PEPCK has mitochondrial and cytosolic isoforms. Oxaloacetate itself does not cross the inner mitochondrial membrane, so it may be transferred indirectly as malate or aspartate.

Fructose-1,6-bisphosphatase hydrolyzes fructose-1,6-bisphosphate to fructose-6-phosphate. Glucose-6-phosphatase hydrolyzes glucose-6-phosphate to free glucose, principally in liver and kidney and to a lesser extent in small intestine. These hydrolysis reactions do not directly consume ATP.

These bypasses are illustrated on the figure on the following page.

\begin{center}\includegraphics[width=0.82\textwidth,height=0.68\textheight,keepaspectratio]{\reviewnotesimagepath d4-figure-gluconeogenesis-bypasses.png}\par\captionof{figure}{Enzymatic bypasses of the irreversible steps of glycolysis during gluconeogenesis.}\labfig{appendix-figure-d4-gluconeogenesis-bypasses}\end{center}

\begin{center}\includegraphics[width=0.72\textwidth,height=0.85\textheight,keepaspectratio]{\reviewnotesimagepath d4-figure-gluconeogenesis-bypass-comparison.png}\par\captionof{figure}{Comparison of glycolysis and gluconeogenesis, including the bypass reactions and enzymes used in gluconeogenesis.}\labfig{appendix-figure-d4-gluconeogenesis-bypass-comparison}\end{center}

Phosphoenolpyruvate may also be generated by means of an alternate pathway (the pathway chosen depends on the availability of various precursors and NADH). The alternate pathway involves the generation of oxaloacetate, and conversion into malate (oxaloacetate has no mitochondrial transporters) by mitochondrial malate dehydrogenase, which is then converted by cytosolic malate dehydrogenase into oxaloacetate and subsequently PEP by cytosolic PEP carboxykinase. This process increases cytosolic NADH, and hence is a pathway which activates when NADH is depleted.

When cytosolic NADH is available, lactate can be oxidized to pyruvate by lactate dehydrogenase, generating NADH; pyruvate can then enter gluconeogenesis. The direction of this reversible reaction depends on substrate and redox conditions.

\begin{center}\includegraphics[width=0.82\textwidth,height=0.68\textheight,keepaspectratio]{\reviewnotesimagepath d4-figure-alternatively-nadh-high-cytosol-utilized.png}\par\captionof{figure}{Lactate oxidation supplies pyruvate for gluconeogenesis and generates cytosolic NADH.}\labfig{appendix-figure-d4-alternatively-nadh-high-cytosol-utilized-convert}\end{center}

Gluconeogenesis may also produce glucose from substances other than pyruvate, including 4, 5 or 6 carbon intermediates from the TCA cycle. These intermediates may in turn be generated from various amino acids found in the cell. These amino acids are referred to as the glucogenic amino acids. The pathways of glucose generation through the breakdown of amino acids is shown below.

\begin{center}\includegraphics[width=0.82\textwidth,height=0.68\textheight,keepaspectratio]{\reviewnotesimagepath d4-figure-gluconeogenesis-produce-glucose-substances-pyruvate.png}\par\captionof{figure}{Lactate, glycerol, and glucogenic amino acids provide carbon for gluconeogenesis.}\labfig{appendix-figure-d4-gluconeogenesis-produce-glucose-substances-pyruvate-carbon}\end{center}

Primary means of transport for amino groups originating from extrahepatic tissue. After the removal of amino groups, they can then undergo gluconeogenesis.

\subsection{The Cori Cycle}

The Cori cycle transfers lactate produced by glycolysis in exercising muscle and red blood cells to the liver, where it can be converted to glucose by gluconeogenesis. Glucose may then return to peripheral tissues. The cycle supports lactate recycling and hepatic glucose output; muscle fatigue and cramps are multifactorial and should not be attributed simply to lactate accumulation.

\begin{center}\includegraphics[width=0.82\textwidth,height=0.68\textheight,keepaspectratio]{\reviewnotesimagepath d4-figure-cori-cycle-regenerates-glucose-lactate.png}\par\captionof{figure}{The Cori cycle transfers lactate from peripheral tissues to the liver for glucose production.}\labfig{appendix-figure-d4-cori-cycle-regenerates-glucose-lactate-byproduct}\end{center}

\subsection{Fates of Lactate}

Monocarboxylate transporters move lactate together with a proton across cell membranes. Transport is reversible and depends on concentration gradients; lactate may be oxidized to pyruvate in tissues with appropriate metabolic conditions.

\begin{center}\includegraphics[width=0.82\textwidth,height=0.68\textheight,keepaspectratio]{\reviewnotesimagepath d4-figure-export-lactate-muscle-requires-transporter.png}\par\captionof{figure}{Monocarboxylate transporters mediate lactate movement across cell membranes.}\labfig{appendix-figure-d4-export-lactate-muscle-requires-transporter-monocarboxylate}\end{center}

Many tumors show increased glucose uptake and glycolytic lactate production, sometimes despite adequate oxygen availability (aerobic glycolysis). Hypoxia can further stabilize hypoxia-inducible factors and increase expression of glycolytic enzymes and glucose transporters. Mitochondrial oxidation is not universally lost in cancer; metabolic phenotypes vary across and within tumors.

Glycolysis and pentose-phosphate enzymes are studied as potential therapeutic targets, but experimental compounds should not be presented as established chemotherapy. Inhibiting these pathways may affect tumor and normal cells, and proposed treatments require evidence of safety and clinical benefit.

\subsection{The Pentose Phosphate Pathway}

The pentose phosphate pathway is a metabolic pathway which is parallel to glycolysis. It generates NADPH and pentoses (5 carbon molecules). The pentose phosphate pathway oxidizes glucose-6-phosphate into pentose phosphates (from 6-carbon to 5-carbon compounds). Notably, pentose phosphates are a key component in nucleic acids.

The pentose phosphate pathway uses NADP\textsuperscript{+} as an electron acceptor in its oxidative phase, generating NADPH and ribulose-5-phosphate. Ribulose-5-phosphate can be converted to ribose-5-phosphate, a precursor for nucleotide synthesis. The pathway has oxidative and non-oxidative phases.

\begin{center}\includegraphics[width=0.82\textwidth,height=0.68\textheight,keepaspectratio]{\reviewnotesimagepath d4-figure-pentose-phosphate-pathway-generates-nadp.jpg}\par\captionof{figure}{The oxidative pentose phosphate pathway generates NADPH and ribulose-5-phosphate.}\labfig{appendix-figure-d4-pentose-phosphate-pathway-generates-nadp-electron}\end{center}

In the first (oxidative) phase of the pentose phosphate pathway, the oxidation of glucose-6-phosphate by the enzyme \textbf{glucose-6-phosphate dehydrogenase} occurs. One proton and electron is removed. The electron is subsequently stored in a NADPH molecule. This results in the product 6-phosphoglucono-\(\delta\)-lactone. Following this, lactonase turns 6-phosphoglucono-\(\delta\)-lactone into 6-phosphogluconate, while also generating a hydrogen ion from water. Next, 6-phosphogluconate dehydrogenase converts 6-phosphogluconate to ribulose-5-phosphate by removing a proton and electron (to NADPH), and generating one molecule of CO2 (by oxidation and decarboxylation). Finally, ribulose-5-phosphate is then isomerized by phosphopentose isomerase into ribose-5-phosphate.

*Note that NADPH prevents oxidative damage.

\begin{center}\includegraphics[width=0.82\textwidth,height=0.68\textheight,keepaspectratio]{\reviewnotesimagepath d4-figure-nadph-prevents-oxidative-damage.png}\par\captionof{figure}{Oxidative pentose phosphate pathway reactions and NADPH production}\labfig{appendix-figure-d4-pentose-phosphate-oxidative-phase}\end{center}

\begin{center}\includegraphics[width=0.82\textwidth,height=0.68\textheight,keepaspectratio]{\reviewnotesimagepath d4-figure-nadph-prevents-oxidative-damage-02.png}\par\captionof{figure}{Structural formula of a pentose phosphate pathway intermediate}\labfig{appendix-figure-d4-pentose-phosphate-intermediate}\end{center}

Glucose-6-phosphate dehydrogenase (G6PD) deficiency is an inherited disorder in which reduced enzyme activity limits NADPH production in red blood cells. This impairs protection against oxidative stress and can precipitate hemolysis after certain infections, medicines, or foods. Severity varies with the underlying variant.

\begin{center}\includegraphics[width=0.82\textwidth,height=0.68\textheight,keepaspectratio]{\reviewnotesimagepath d4-figure-glucose-phosphate-dehydrogenase-deficiency-g6pd.jpg}\par\captionof{figure}{Reduced G6PD activity limits NADPH generation and increases susceptibility to oxidative stress.}\labfig{appendix-figure-d4-glucose-phosphate-dehydrogenase-deficiency-g6pd-genetic}\end{center}

The non-oxidative phase reversibly rearranges pentose phosphates into glycolytic intermediates, including fructose-6-phosphate and glyceraldehyde-3-phosphate. These reactions do not directly generate NADPH; their direction depends on cellular needs for ribose-5-phosphate and glycolytic intermediates.

\begin{center}\includegraphics[width=0.82\textwidth,height=0.68\textheight,keepaspectratio]{\reviewnotesimagepath d4-figure-pentose-phosphate-pathway-continues-non.png}\par\captionof{figure}{Reversible non-oxidative pentose-phosphate reactions interconvert sugar phosphates.}\labfig{appendix-figure-d4-pentose-phosphate-pathway-continues-non-oxidative-phase}\end{center}

Glucose-6-phosphate can enter either glycolysis or the pentose phosphate pathway. Flux through the pentose phosphate pathway rises when NADP\textsuperscript{+} is available and the cell requires NADPH or ribose-5-phosphate; NADPH inhibits glucose-6-phosphate dehydrogenase by feedback regulation.

\begin{center}\includegraphics[width=0.82\textwidth,height=0.68\textheight,keepaspectratio]{\reviewnotesimagepath d4-figure-relation-glycolysis-molecule-reactions-glucose.png}\par\captionof{figure}{Glycolytic intermediates connect glucose metabolism to other biosynthetic pathways.}\labfig{appendix-figure-d4-relation-glycolysis-molecule-reactions-glucose-phosphate}\end{center}

\subsection{Glycogen}

Glycogen is stored as granules, primarily in liver and skeletal muscle. During glycogenolysis, glycogen phosphorylase releases glucose-1-phosphate from nonreducing ends; a debranching enzyme handles the alpha-1,6 branch points. The liver can convert glucose-6-phosphate to free glucose for release, whereas skeletal muscle uses glucose-6-phosphate locally.

\begin{center}\includegraphics[width=0.82\textwidth,height=0.68\textheight,keepaspectratio]{\reviewnotesimagepath d4-figure-glycogen-found-muscle-liver-tissues.png}\par\captionof{figure}{Glycogen granules are stored mainly in liver and skeletal muscle.}\labfig{appendix-figure-d4-glycogen-found-muscle-liver-tissues-granules}\end{center}

Glycogenesis stores glucose as glycogen, a branched polymer with alpha-1,4 linkages and alpha-1,6 branch points. UDP-glucose is formed by UDP-glucose pyrophosphorylase and added to a glycogen primer by glycogen synthase, which extends alpha-1,4 linkages; the branching enzyme forms alpha-1,6 linkages.

\begin{center}\includegraphics[width=0.82\textwidth,height=0.68\textheight,keepaspectratio]{\reviewnotesimagepath d4-figure-inversely-synthesis-glycogen-glucose-storage.png}\par\captionof{figure}{Reciprocal regulation of glycogen synthesis and breakdown.}\labfig{appendix-figure-d4-inversely-synthesis-glycogen-glucose-storage-glycogenesis}\end{center}

\begin{center}\includegraphics[width=0.82\textwidth,height=0.68\textheight,keepaspectratio]{\reviewnotesimagepath d4-figure-inversely-synthesis-glycogen-glucose-storage-02.png}\par\captionof{figure}{Reciprocal regulation of glycogen synthesis and breakdown.}\labfig{appendix-figure-d4-inversely-synthesis-glycogen-glucose-storage-glycogenesis-02}\end{center}

Glycogen synthesis is promoted by insulin, which favors glycogen-synthase activity and inhibits glycogen breakdown. In skeletal muscle, insulin also promotes GLUT4 translocation and glucose uptake. These effects do not increase glucagon production.

\begin{center}\includegraphics[width=0.82\textwidth,height=0.68\textheight,keepaspectratio]{\reviewnotesimagepath d4-figure-synthesis-glycogen-insulin-dependant-binds.png}\par\captionof{figure}{Insulin signaling promotes glucose uptake and glycogen synthesis in skeletal muscle.}\labfig{appendix-figure-d4-synthesis-glycogen-insulin-dependant-binds-muscle}\end{center}

\subsection{Blood Glucose Homeostasis}

The blood glucose homeostasis refers to the control of plasma blood glucose level. This homeostasis balances between the amount of absorbed glucose and the amount used by cells of the body.

\begin{center}\includegraphics[width=0.82\textwidth,height=0.68\textheight,keepaspectratio]{\reviewnotesimagepath d4-figure-blood-glucose-homeostasis-control-plasma.png}\par\captionof{figure}{Hormonal regulation of plasma glucose by insulin and counter-regulatory hormones.}\labfig{appendix-figure-d4-blood-glucose-homeostasis-control-plasma-level}\end{center}

Blood-glucose reference ranges depend on fasting status, timing, and clinical context; diabetes is diagnosed using validated laboratory criteria rather than a single informal “optimal” range. Low glucose can impair neurological function, with thresholds depending on clinical context. Hemoglobin A1c reflects average glycemia over roughly the preceding two to three months and is one accepted test for diabetes when interpreted using diagnostic criteria.

Blood glucose is regulated primarily by insulin from pancreatic beta cells and glucagon from alpha cells. Insulin promotes glucose uptake in insulin-responsive tissues and favors storage, while glucagon increases hepatic glucose output through glycogenolysis and gluconeogenesis.

\begin{center}\includegraphics[width=0.82\textwidth,height=0.68\textheight,keepaspectratio]{\reviewnotesimagepath d4-figure-control-blood-glucose-level-achieved.png}\par\captionof{figure}{Insulin and glucagon coordinate blood-glucose homeostasis.}\labfig{appendix-figure-d4-control-blood-glucose-level-achieved-primarily}\end{center}

Glycogen is not evenly distributed around the body. It is stored in specific locations. The primary storage of glycogen for release into the bloodstream is in the liver (3-5\% of the tissue mass 75g/70kg). Even though muscles also store glycogen (0.5-1.0\% of tissue mass 250g/70kg), they are not able to release it into the bloodstream. The graph below showcases the source of glucose at various points during the typical 24 hour day. Notice that glucagon acts as the main source of glucose during the day between meals, and gluconeogenesis increases after dinner to replenish the glucose supply.

\begin{center}\includegraphics[width=0.82\textwidth,height=0.68\textheight,keepaspectratio]{\reviewnotesimagepath d4-figure-glycogen-evenly-distributed-around-body.png}\par\captionof{figure}{Tissue distribution of glycogen stores and their metabolic roles.}\labfig{appendix-figure-d4-glycogen-evenly-distributed-around-body-stored}\end{center}

Shown below are the major links between carbohydrates, fats and proteins. Namely, fats in the form of triglycerides may be hydrolyzed into glycerol and fatty acids. These breakdown products may enter into the glycolytic cycle and be metabolized, with glycerol being transferred into glyceraldehyde-3-phosphate, fatty acids to acetyl coenzyme A. Meanwhile, proteins which are hydrolyzed into amino acids will have their amino groups removed before entering the krebs cycle as acetyl CoA or the citric acid cycle.

\begin{center}\includegraphics[width=0.82\textwidth,height=0.68\textheight,keepaspectratio]{\reviewnotesimagepath d4-figure-major-links-between-carbohydrates-fats.png}\par\captionof{figure}{Shared intermediates link carbohydrate, lipid, and protein metabolism.}\labfig{appendix-figure-d4-major-links-between-carbohydrates-fats-proteins}\end{center}

\section{Lipid Metabolism}

\subsection{Lipids}

Lipids are organic compounds which contain fat or fat like elements inside of their structure. Examples include fat, oils, phospholipids and cholesterol. They are hydrophobic and nonpolar, and mainly supply energy for cellular processes and act as structural components inside of cell membranes (eg. phospholipids). The most common dietary lipids are triglycerides, a type of fat. Excess amounts lead to obesity and related disorders, while deficiency leads to weight loss and hormonal imbalances. As such, various lipids in serum (eg. HDL, LDL) may be used as indicators of health.

Triglycerides are composed of one glycerol and three fatty acid molecules bound by ester bonds formed during dehydration. Note that in water, the carboxyl groups may lose their hydrogen and become negatively charged. Typical conformations for triglycerides and glycerol and fatty acids are shown below. Additionally, note that glycerol may also be used in the synthesis of glucose.

\begin{center}\includegraphics[width=0.82\textwidth,height=0.68\textheight,keepaspectratio]{\reviewnotesimagepath d4-figure-triglycerides-composed-glycerol-three-fatty.png}\par\captionof{figure}{A triacylglycerol consists of glycerol esterified to three fatty acids.}\labfig{appendix-figure-d4-triglycerides-composed-glycerol-three-fatty-acid}\end{center}

Triglycerides may be classified based on their degree of saturation, which refers to the number of hydrogen atoms in a given lipid. Saturated fats are fats which have the maximum possible number of hydrogen atoms (meaning that all carbons have single bonds), while unsaturated fats have double bonds, which leads to them not having the full number of hydrogens. In nature, saturated fats are most common in plant oils and nature in general.

Saturated and unsaturated fats have different characteristics owing to their structure and intermolecular interactions. Saturated fat contains molecules which are densely packed together, leading to many intermolecular interactions which means that they are solid at room temperature. On the other hand, unsaturated fats have double bonds, which leads to kinks in the triglyceride structure, which In turn causes a decreased packing effect and liquids at room temperature.

\subsection{The Function of Lipids}

Different lipid structures each express different functions. There are three main functions of lipids in the body. Firstly, lipids in adipose tissue (usually referred to as storage lipids) are used to store energy in the form of triglycerides (they also act as structural components). Enzymes in adipose tissue may oxidise fat and release it into circulation for other cells to use as energy.

Phospholipids (which are aliphatic molecules) act as a part of double layer biological membranes, playing a core role in controlling cell structure and permeability by stopping essential water soluble substances into and out of the cell. This plays a part in the compartmentalization of cells. There are 5 main types of membrane lipids with different structures: glycerophospholipids, galactolipids/sulfolipids, archaeal tetraether lipids, sphingolipids and sterols.

Both storage and membrane lipids act in passive roles. However, some lipids (which are present in small amounts) play an active role in the body, being signals, cofactors and pigments. For example, Cholesterol acts as a hormonal synthesis precursor.

The differentiating factor between the active and passive roles of lipids is whether they function in one location or are spread throughout the body. If a molecule needs to be transported, it is considered to have an active role.

\begin{center}\includegraphics[width=0.82\textwidth,height=0.68\textheight,keepaspectratio]{\reviewnotesimagepath d4-figure-between-active-passive-lipids-whether.png}\par\captionof{figure}{Structural and signaling roles of biologically active lipids.}\labfig{appendix-figure-d4-between-active-passive-lipids-whether-function}\end{center}

Storage of energy in the form of fat is done both by plants and animals in specific cell types. In animals, fat is stored in adipose tissue as triglycerides or triacylglycerols (in fat droplets). This fat acts as a supply of energy for the organism.

The different structure of lipids leads to different functions; shown below are the typical membrane lipids.


Sterols differ from other lipids because they are composed of 4 filled links of steroid nucleus, which is their major structure. Sterols lead to the production of hormones.

\begin{center}\includegraphics[width=0.82\textwidth,height=0.68\textheight,keepaspectratio]{\reviewnotesimagepath d4-figure-sterols-differ-lipids-composed-filled.png}\par\captionof{figure}{The fused-ring sterol structure distinguishes sterols from other lipid classes.}\labfig{appendix-figure-d4-sterols-differ-lipids-because-composed-filled}\end{center}

Finally, lipids in active roles act as signal molecules, and are present in smaller amounts. They have an active role in the metabolic process either as messengers (such as hormones) or metabolites. Signal molecules are often carried in the bloodstream from one tissue to another. One example is arachidonate, which is a precursor for the formation of eicosanoids. eicosanoids are in turn the precursor of prostaglandins, thromboxane and leukotrienes (paracrine hormones).

\subsection{Dietary Lipid Processing}

\begin{center}\includegraphics[width=0.82\textwidth,height=0.68\textheight,keepaspectratio]{\reviewnotesimagepath d4-figure-dietary-lipid-processing.png}\par\captionof{figure}{Emulsification, digestion, and absorption of dietary lipids.}\labfig{appendix-figure-d4-dietary-lipid-processing}\end{center}

Pancreatic lipases digest dietary triacylglycerols into fatty acids and monoacylglycerols. Long-chain products enter enterocytes from micelles, are re-esterified to triacylglycerols, and are packaged with apolipoprotein B-48 into chylomicrons, which enter lymph before reaching the blood. Glycerol and many short- and medium-chain fatty acids enter portal blood and circulate largely bound to albumin; chain-length categories vary by convention.

Inside enterocytes, long-chain fatty acids are activated to fatty acyl-CoA and re-esterified with monoacylglycerols to form triacylglycerols. The triacylglycerols are packaged into chylomicrons, which transport dietary lipids through lymph and blood; chylomicrons do not facilitate uptake across the intestinal cell membrane.

\begin{center}\includegraphics[width=0.82\textwidth,height=0.68\textheight,keepaspectratio]{\reviewnotesimagepath d4-figure-larger-molecules-diacylglycerol-monoacylglycerol-long.png}\par\captionof{figure}{Intestinal digestion converts dietary triacylglycerols into absorbable lipid products.}\labfig{appendix-figure-d4-larger-molecules-diacylglycerol-monoacylglycerol-long-chain}\end{center}

Chylomicrons are lipoproteins assembled by enterocytes to transport dietary lipids. The liver secretes very-low-density lipoprotein (VLDL) to transport endogenous triacylglycerols; VLDL is remodeled in the circulation to intermediate-density lipoprotein (IDL) and then low-density lipoprotein (LDL).

Lipoproteins are composed of phospholipids in a monolayer, with hydrophilic phosphate groups on the outside and a hydrophobic interior containing cholesterol, esters and triglycerides. The phospholipid layer integrates cholesterol to help stabilize the structure, as well as apo(lipo)proteins.

\begin{center}\includegraphics[width=0.82\textwidth,height=0.68\textheight,keepaspectratio]{\reviewnotesimagepath d4-figure-lipoproteins-composed-phospholipids-monolayer-hydrophilic.png}\par\captionof{figure}{Lipoprotein structure places a phospholipid monolayer around a hydrophobic core.}\labfig{appendix-figure-d4-lipoproteins-composed-phospholipids-monolayer-hydrophilic-phosphate}\end{center}

Table~\ref{tab:appendix-table-d4-lipoprotein-classes} compares selected physical and compositional properties of the major lipoprotein classes.
\pagelayout{wide}
\begin{center}
\captionof{table}{Selected properties of plasma lipoprotein classes.}\labtab{appendix-table-d4-lipoprotein-classes}
\small
\renewcommand{\arraystretch}{1.15}
\begin{tabularx}{\textwidth}{p{0.15\textwidth}p{0.17\textwidth}p{0.17\textwidth}Xp{0.15\textwidth}}
\toprule
\textbf{Particle} & \textbf{Density (kg/L)} & \textbf{Main component} & \textbf{Apolipoproteins\textsuperscript{*}} & \textbf{Diameter (nm)} \\
\midrule
Chylomicrons & 0.95 & TG & B48 (A, C, E) & 75--1200 \\

VLDL & 0.95--1.006 & TG & B100 (A, C, E) & 30--80 \\

IDL & 1.006--1.019 & TG \& cholesterol & B100, E & 25--35 \\

LDL & 1.019--1.063 & Cholesterol & B100 & 18--25 \\

HDL & 1.063--1.210 & Protein & AI, AII (C, E) & 5--12 \\
\bottomrule
\end{tabularx}
\par\smallskip
\begin{minipage}{\textwidth}
\footnotesize\raggedright
TG, triacylglycerol (triglyceride); VLDL, very-low-density lipoprotein; IDL, intermediate-density lipoprotein; HDL, high-density lipoprotein. On electrophoresis, VLDL are called pre-\(\beta\)-lipoproteins, LDL \(\beta\)-lipoproteins, and HDL \(\alpha\)-lipoproteins. The asterisk appears in the source heading without an accompanying definition.
\end{minipage}
\end{center}
\pagelayout{margin}

Lipoprotein classes differ in density, size, lipid composition, and apolipoproteins. VLDL is secreted by the liver and progressively loses triacylglycerol to form IDL and LDL. LDL delivers cholesterol to tissues through receptor-mediated uptake; elevated LDL cholesterol is associated with atherosclerotic cardiovascular risk. HDL participates in reverse cholesterol transport to the liver, where cholesterol may be recycled or converted to bile acids and excreted.

Apolipoproteins provide structural support and act as enzyme cofactors or receptor ligands. ApoB-48 is characteristic of chylomicrons, whereas ApoB-100 is present in VLDL and its IDL and LDL descendants. ApoC-II activates lipoprotein lipase, and ApoE contributes to hepatic remnant uptake.

Lipoprotein lipase, anchored to capillary endothelium and activated by ApoC-II, hydrolyzes triacylglycerols in chylomicrons and VLDL, releasing fatty acids for tissue uptake. Chylomicron remnants and some IDL are taken up by the liver through apolipoprotein-dependent receptors; LDL particles are cleared primarily through LDL receptors recognizing ApoB-100.

The lipid transportation pathway in the human body starts from chylomicrons and VLDL. It is often referred to as the fuel transport pathway. This pathway is closely related to the fed and fast cycle; during the fed stage and food absorption, chylomicron transports triglycerides to peripheral cells and tissues, where lipoprotein lipase and receptors bind and extract triglycerides and fatty acid into the cell. The remnants of this process are transported back to the liver and metabolized. Remnants bind to a receptor on the surface of hepatocytes (the LRP receptor), where they are used to resynthesize lipoproteins. The process mentioned above (and the dark line in the diagram) is the exogenous pathway lipid transportation mechanism (as the fat comes from ingestion).

Meanwhile, the endogenous pathway of lipid transportation is highlighted in the pale line. This pathway utilizes VLDL, a factor synthesized by hepatocytes. In the fasting stage or during high glucose (as glucose’s energy can be stored as fat), fat is transported from the liver to peripheral tissues. VLDL remnants are then returned to the liver and taken up by the ApoB/E receptor.

The alternate (overflow) pathway that is activated when the liver is not able to metabolize remnants at a fast enough rate, the LDL particles are generated from the remnants. These LDL particles (which now transport cholesterol) bind to peripheral cells again to a cholesterol (B/E receptor). If there is a large amount of cholesterol transportation, there is a high likelihood of blood vessel plaque formation, and subsequent heart failure or strokes.

\begin{center}\includegraphics[width=0.82\textwidth,height=0.68\textheight,keepaspectratio]{\reviewnotesimagepath d4-figure-alternate-overflow-pathway-activated-liver.png}\par\captionof{figure}{Hepatic ketogenesis increases when fatty-acid oxidation exceeds TCA-cycle capacity.}\labfig{appendix-figure-d4-alternate-overflow-pathway-activated-liver-able}\end{center}

The reverse cholesterol transport pathway is responsible for the transport of HDL (it is linked closely with the fuel transport pathway). HDL is a protein dense carrier which functions to remove cholesterol from peripheral cells. The pathway starts from the synthesis of HDL in the liver. HDL then removes cholesterol from peripheral cells, causing it to become bigger. HDL then binds to the liver scavenger receptor and releases its cholesterol (the direct pathway). During the transportation of cholesterol, HDL may exchange cholesterol contents (in the form of cholesterol esters) for triglycerides with VLDL, LDL or IDL. This indirect pathway reduces plaque formation by absorbing cholesterol (antiartheroschlerotic). The pathway is shown below along with a general overview of the two pathways.

\begin{center}\includegraphics[width=0.82\textwidth,height=0.68\textheight,keepaspectratio]{\reviewnotesimagepath d4-figure-reverse-cholesterol-transport-pathway-responsible.png}\par\captionof{figure}{HDL-mediated reverse cholesterol transport returns peripheral cholesterol to the liver.}\labfig{appendix-figure-d4-reverse-cholesterol-transport-pathway-responsible-hdl}\end{center}

\begin{center}\includegraphics[width=0.82\textwidth,height=0.68\textheight,keepaspectratio]{\reviewnotesimagepath d4-figure-reverse-cholesterol-transport-pathway-responsible-02.png}\par\captionof{figure}{HDL-mediated reverse cholesterol transport returns peripheral cholesterol to the liver.}\labfig{appendix-figure-d4-reverse-cholesterol-transport-pathway-responsible-hdl-02}\end{center}

\subsection{Lipid Catabolism}

Once triglycerides enter into the cell, they are first broken down by lipase (through hydrolysis). Fatty acids are then converted through beta oxidation into acetyl coenzyme A, which then enters into the citric acid cycle (TCA cycle). Alternatively, conversion into ketone bodies is also possible. The glycerol which has been released by hydrolysis may be converted into glucose by the gluconeogenesis pathway. This process is shown below. Note that the oxidation of fatty acids is the main source of energy in the catabolism of lipids (which uses enzymes located in the mitochondrial matrix). Also note that sterols are not catabolized, but excreted.

Glycerol may also be broken down into glyceraldehyde 3 phosphate (the first breakdown is to glycerol-3-phosphate by \textbf{glycerol kinase} using 1 ATP, which is followed by the breakdown into dihydroxyacetone phosphate by glycerol-3-phosphate dehydrogenase), which in turn is converted into pyruvic acid and subsequently acetyl coenzyme A. This is also shown on the following diagrams.

\begin{center}\includegraphics[width=0.82\textwidth,height=0.68\textheight,keepaspectratio]{\reviewnotesimagepath d4-figure-glycerol-broken-down-glyceraldehyde-phosphate.png}\par\captionof{figure}{Glycerol enters carbohydrate metabolism after conversion to dihydroxyacetone phosphate.}\labfig{appendix-figure-d4-glycerol-broken-down-glyceraldehyde-phosphate-breakdown}\end{center}

\begin{center}\includegraphics[width=0.82\textwidth,height=0.68\textheight,keepaspectratio]{\reviewnotesimagepath d4-figure-glycerol-broken-down-glyceraldehyde-phosphate-02.png}\par\captionof{figure}{Glycerol enters carbohydrate metabolism after conversion to dihydroxyacetone phosphate.}\labfig{appendix-figure-d4-glycerol-broken-down-glyceraldehyde-phosphate-breakdown-02}\end{center}

\subsubsection{Fasting Metabolism}

Excess energy can be stored in adipose tissue as triacylglycerols and mobilized during fasting or increased energy demand. In humans, catecholamines are important acute stimulators of adipocyte lipolysis, while insulin suppresses it; the direct effect of glucagon on human adipose tissue is limited and species-dependent. Caffeine may influence catecholaminergic signaling but is not itself a lipolytic hormone.

Catecholamine binding to beta-adrenergic receptors activates Gs, adenylyl cyclase, and cAMP-dependent protein kinase A. Phosphorylation of perilipin and hormone-sensitive lipase contributes to lipid-droplet mobilization; adipose triglyceride lipase, hormone-sensitive lipase, and monoacylglycerol lipase act sequentially to release fatty acids and glycerol. Circulating non-esterified fatty acids travel bound to albumin and can be taken up by tissues.

\begin{center}\includegraphics[width=0.82\textwidth,height=0.68\textheight,keepaspectratio]{\reviewnotesimagepath d4-figure-hormonal-receptor-activates-adipocyte-membrane.png}\par\captionof{figure}{Hormonal signals regulate adipocyte lipolysis and fatty-acid release.}\labfig{appendix-figure-d4-hormonal-receptor-activates-adipocyte-membrane-simultaneously}\end{center}

In the fasting stage, the liver is the primary source of glucose for the brain. The liver produces glucose through various mechanisms, including the breakdown of glycogen and gluconeogenesis. In this process, glycerol (a product of triglyceride lysis) is also used in gluconeogenesis, while free fatty acids are used as the primary source of energy during the production of glucose. Note that during the fasting stage, glycolysis is suppressed to conserve glucose for export out of the liver.

\begin{center}\includegraphics[width=0.82\textwidth,height=0.68\textheight,keepaspectratio]{\reviewnotesimagepath d4-figure-fasting-stage-liver-primary-source.png}\par\captionof{figure}{Hepatic glycogenolysis and gluconeogenesis help maintain blood glucose during fasting.}\labfig{appendix-figure-d4-fasting-stage-liver-primary-source-glucose}\end{center}

\subsubsection{Fatty Acids and the Mitochondria}

Before generating energy, fatty acids are activated into their CoA form (fatty acyl CoA) in the cytosolic side of the outer mitochondrial membrane before transportation into the matrix. The short chain (2-4 FA) and medium chain (4-12 FA) fatty acids are activated inside of the mitochondria, while long chain fatty acids (12-20 FA) are activated in the cytoplasm and subsequently transported by the carnitine cycle.

Long chain fatty acids are activated into their CoA derivatives by fatty acyl CoA synthetase (also called long chain acetyl coenzyme A synthetase) on the outer cytosolic leaflet of the mitochondria. Acyl CoA (the product) is then bound to carnitine (by removing the CoA group). This transforms Acyl CoA into acylcarnitine. Acylcarnitine is then transported by CPT1 (carnitine palmitoyltransferase 1) into the intermembrane space of the mitochondria. At the intermembrane space, CACT on the outwards facing side of the inner mitochondrial membrane transports acylcarnitine into the matrix while also removing free carnitine from the matrix (as it acts as an antiporter).

Upon entry of acylcarnitine into the mitochondrial matrix, CPT2 subsequently regenerates the Coenzyme A group and releases the carnitine shuttle which travels back to the cytoplasm, where the process repeats. Simultaneously, the Acyl-CoA undergoes beta oxidation. An overview of this process is shown on the following diagram.

\begin{center}\includegraphics[width=0.82\textwidth,height=0.68\textheight,keepaspectratio]{\reviewnotesimagepath d4-figure-upon-entry-acylcarnitine-mitochondrial-matrix.png}\par\captionof{figure}{The carnitine shuttle transports long-chain fatty acyl groups into mitochondria.}\labfig{appendix-figure-d4-upon-entry-acylcarnitine-mitochondrial-matrix-cpt2}\end{center}

\begin{center}\includegraphics[width=0.82\textwidth,height=0.68\textheight,keepaspectratio]{\reviewnotesimagepath d4-figure-upon-entry-acylcarnitine-mitochondrial-matrix-02.png}\par\captionof{figure}{The carnitine shuttle transports long-chain fatty acyl groups into mitochondria.}\labfig{appendix-figure-d4-upon-entry-acylcarnitine-mitochondrial-matrix-cpt2-02}\end{center}

\subsubsection{Fatty Acid Oxidation}

Fatty acid oxidation consists of three main stages: beta oxidation, where long chain fatty acids are oxidised into acetyl CoA residue, the citric acid cycle, and then finally the electron transport chain. In all, this process results in 108 ATP for palmitoyl-CoA.

\begin{center}\includegraphics[width=0.82\textwidth,height=0.68\textheight,keepaspectratio]{\reviewnotesimagepath d4-figure-fatty-acid-oxidation-consists-three.png}\par\captionof{figure}{Each beta-oxidation cycle shortens a fatty acyl-CoA and generates reducing equivalents.}\labfig{appendix-figure-d4-fatty-acid-oxidation-consists-three-stages}\end{center}

Beta oxidation is a 4 step cycle which starts with the dehydrogenation of the fatty acid by \textbf{acyl-CoA dehydrogenase}. Beta oxidation starts from the beta (second) carbon of each molecule. One beta oxidation results in 1 NADH and 1 FADH\textsubscript{2}, and it is repeated every 2 carbons for the whole molecule. which starts between the alpha and beta carbon. Acyl-CoA dehydrogenase deprotonates (removes a hydrogen from) the alpha carbon and places it into FADH\textsubscript{2}. This leads to a double bond in the molecule (which is converted to Enoyl-CoA). Following the deprotonation, enoyl CoA hydrates the double bond, resulting in the production of a hydroxyl group on the beta carbon molecule. Dehydrogenase then removes another hydrogen into NADH (and removes another to a free proton), which results in ketoacyl-CoA. This is then cleaved by thiolase (also known as acyl-CoA acetyltransferase) into acetyl-CoA (the product) and the remaining carbon structure. This process is shown below. Note that it repeats until the whole molecule becomes a two-carbon molecule. For instance, a 14 carbon atom undergoes this cycle 6 times until the remaining product is two acyl-CoA. This cycle is shown on the right.

\begin{center}\includegraphics[width=0.82\textwidth,height=0.68\textheight,keepaspectratio]{\reviewnotesimagepath d4-figure-beta-oxidation-step-cycle-starts.png}\par\captionof{figure}{The beta-oxidation cycle produces acetyl-CoA, NADH, and FADH2.}\labfig{appendix-figure-d4-beta-oxidation-step-cycle-starts-dehydrogenation}\end{center}

For unsaturated fatty acids (partially oxidised with double bounds) will cause the decreased generation of FADH\textsubscript{2}, as the deprotonation step is bypassed for the carbon double bond linkage. This leads to a lower quantity of ATP being generated. Polyunsaturated fatty acids will have a cis configuration at 3 carbon intervals, this results in the need to adapt intermediates into a trans geometry (achieved by isomerase). After this, the reaction will proceed in 2 carbon steps as usual. Note that this reaction requires additional isomerase and oxidoreductase enzymes to form two carbon compounds. This process is shown below for unsaturated fatty acids. For polyunsaturated fatty acids, this process is simply repeated for as many times as there are double bonds.

\begin{center}\includegraphics[width=0.82\textwidth,height=0.68\textheight,keepaspectratio]{\reviewnotesimagepath d4-figure-unsaturated-fatty-acids-partially-oxidised.png}\par\captionof{figure}{Additional enzymes enable beta-oxidation of unsaturated fatty acids.}\labfig{appendix-figure-d4-unsaturated-fatty-acids-partially-oxidised-double}\end{center}

For odd chain fatty acids, at the end of beta oxidation, propionyl-CoA (a 3-carbon molecule) will be formed. This is converted through a process with multiple steps (which involves 3 enzymes and biotin/cobalamin cofactors). Firstly, propionyl-CoA carboxylase changes the structure of propionyl-CoA to succinyl-CoA to enter the citric acid cycle. This process is shown below.

\begin{center}\includegraphics[width=0.82\textwidth,height=0.68\textheight,keepaspectratio]{\reviewnotesimagepath d4-figure-odd-chain-fatty-acids-end.png}\par\captionof{figure}{Odd-chain fatty-acid oxidation yields propionyl-CoA, which is converted to succinyl-CoA.}\labfig{appendix-figure-d4-odd-chain-fatty-acids-end-beta}\end{center}

\subsubsection{Peroxisomal β-Oxidation}

Peroxisomes handle the oxidation of very long chain fatty acids. Beta oxidation inside peroxisomes has differences compared to mitochondrial beta oxidation, but they both utilize carnitine shuttles for the transport of fatty acids into the organelle. The peroxisomal system differs in that the acyl-CoA dehydrogenase is substituted by an oxidase, and the production of FADH\textsubscript{2} (along with other type of oxidation) will have the end product of the production of hydrogen peroxide (as FADH\textsubscript{2} in the peroxisome is utilized to manufacture hydrogen peroxide). Furthermore, instead of the usage of NADH inside of the molecule, it is exported to be reoxidised in the cytosol (along with the Acetyl-CoA, which is also exported for use in biosynthesis).

\begin{center}\includegraphics[width=0.82\textwidth,height=0.68\textheight,keepaspectratio]{\reviewnotesimagepath d4-figure-peroxisomes-handle-oxidation-long-chain.png}\par\captionof{figure}{Peroxisomes initiate oxidation of very-long-chain fatty acids.}\labfig{appendix-figure-d4-peroxisomes-handle-oxidation-long-chain-fatty}\end{center}

\subsection{Lipid Anabolism}

https://doctorlib.org/medical/biochemistry/18.html

\subsubsection{Fatty Acid Synthesis}

Lipid anabolism is also referred to as lipogenesis. It occurs in response to excess energy intake as signaled by insulin. To produce lipids, fatty acid (manufactured from acetyl-CoA) and glycerol are required precursors. These precursors are able to produce triglycerides, lipoproteins, non-essential fatty acids, cholesterol and phospholipids. Typically, lipogenesis occurs in the cytosol of both the liver and adipose tissue. The liver mostly acts as the main lipogenesis site, while the adipose tissue acts mostly as storage. Lipogenesis is not the reverse process of fatty acid oxidation, as it is a different process entirely.

Since lipogenesis occurs in the cytosol, there is a need to transport acetyl-CoA from the mitochondria to the cytosol. This is achieved by the use of malate shuttles (also known as the malate-citrate shuttle). The process starts by the conversion of acetyl-CoA into citrate by the addition of oxaloacetate. Citrate transporters (which are also malate antiporters) then move citrate out of the mitochondrial membrane. Upon introduction to the cytosol, ATP citrate lyase uses ATP to cleave citrate back into acetyl-CoA and oxaloacetate. The liberated acetyl-CoA is then used to produce malonyl-CoA, which is a fatty acid precursor. Meanwhile, oxaloacetate is turned back into malate, which is imported into the mitochondria by the malate-citrate antiporter. Malate then changes back to oxaloacetate through the citric acid cycle in the mitochondria. As such, the cycle repeats (if necessary). See the diagram below.

\begin{center}\includegraphics[width=0.82\textwidth,height=0.68\textheight,keepaspectratio]{\reviewnotesimagepath d4-figure-since-lipogenesis-cytosol-there-transport.png}\par\captionof{figure}{Citrate export supplies cytosolic acetyl-CoA for fatty-acid synthesis.}\labfig{appendix-figure-d4-since-lipogenesis-cytosol-there-transport-acetyl-co}\end{center}

The synthesis of fatty acids is a two stage process, with both stages requiring input of acetyl-CoA, and employing multifunctional proteins in a multienzyme complex. In the first stage, malonyl-CoA (a key precursor for elongation) is formed. In the second stage, fatty acid synthase elongates the chain of fatty acids in two carbon increments.

In the first stage of the fatty acid biosynthesis, most of the acetyl-CoA originates from the carbohydrate metabolism. This acetyl-CoA is converted into malonyl-CoA by the use of \textbf{acetyl-CoA carboxylase}, which is covalently attached to biotin. The process begins with carboxylation of the biotin and bicarbonate followed by covalently attaching biotin to the enzyme. This requires energy in the form of ATP. Following this process, ATP is also used to attach the carbon dioxide to acetyl-CoA, producing malonyl-CoA, a 3 carbon molecule. This process can be activated by polymerization of protomers in the presence of citrate and absence of long chain acyl-CoA (which suppresses the polymerization).

Another method of regulating acetyl-CoA carboxylase is by hormonal pathways, namely insulin induced dephosphorylation through protein phosphatase and epinephrine or glucagon based phosphorylation through protein kinase A (phosphate groups deactivate acetyl-CoA carboxylase).

\begin{center}\includegraphics[width=0.82\textwidth,height=0.68\textheight,keepaspectratio]{\reviewnotesimagepath d4-figure-method-regulating-acetyl-co-carboxylase.png}\par\captionof{figure}{Hormonal and allosteric signals regulate acetyl-CoA carboxylase and lipogenesis.}\labfig{appendix-figure-d4-method-regulating-acetyl-co-carboxylase-hormonal-pathways}\end{center}


The second stage of fatty acid synthesis is the fatty acid elongation, which utilizes fatty acid synthase, a multienzyme complex. It is composed of 7 subunit enzymes per polypeptide, and has two polypeptides attached together into a dimer. The structure is shown on the following diagram.

\begin{center}\includegraphics[width=0.82\textwidth,height=0.68\textheight,keepaspectratio]{\reviewnotesimagepath d4-figure-stage-fatty-acid-synthesis-elongation.png}\par\captionof{figure}{Fatty-acid synthesis uses acetyl-CoA, malonyl-CoA, and reducing power to extend the chain.}\labfig{appendix-figure-d4-stage-fatty-acid-synthesis-elongation-utilizes}\end{center}

Each subunit has its own function inside of the polypeptide. The acyl carrier protein (ACP) has a cysteine binding point. It is a highly conserved protein with the function of replacing coenzyme A and binds to the elongating fatty acid chain starting with the binding of malonyl-CoA. The pantetheine serves as a binding site for acetyl-CoA and malonyl-CoA along with cysteine.

The synthesis reaction starts with the binding of acetyl-CoA to the enzyme at the cysteine group and is followed by the binding of malonyl-CoA to a pantetheine group. The molecule then undergoes synthesis, as the acetyl-CoA is incorporated with malonyl-CoA. One carbon dioxide is removed and reduction (the removal of one proton) occurs. Water is then removed and the product (a four carbon molecule) is formed. After the process is complete, another malonyl-CoA group is added and the process repeats, ending with the formation of up to 16 carbon long palmitate from 6 more repetitions of the cycle (see following). The ending point is when thioesterase releases palmitate from the enzyme complex.

\textbf{Need to know: Fatty acid synthase is a multienzyme complex which uses acetyl and malonyl-CoA as a precursor for synthesis. Every cycle, one carbon is removed and malonyl-CoA is added, leading to a total of 2 carbon added per cycle.}

\begin{center}\includegraphics[width=0.82\textwidth,height=0.68\textheight,keepaspectratio]{\reviewnotesimagepath d4-figure-know-fatty-acid-synthase-multienzyme.png}\par\captionof{figure}{Organization of the fatty-acid synthase multienzyme complex.}\labfig{appendix-figure-d4-know-fatty-acid-synthase-multienzyme-complex}\end{center}

\subsubsection{Fatty Acid Elongation}

Fatty acid elongation is the process which extends a fatty acid beyond 16 carbons. It requires an enzyme, \textbf{fatty acid elongase}. This elongation adds ten or more carbon molecules by using 10 carbon (or longer) chains as substrates. This process is the reversal of beta oxidation. The reactions occurring during chain elongation are similar to those involved in fatty acid synthesis, except that the fatty acid is attached to CoA, rather than to the ACP.

\begin{center}\includegraphics[width=0.82\textwidth,height=0.68\textheight,keepaspectratio]{\reviewnotesimagepath d4-figure-fatty-acid-elongation-extends-beyond.png}\par\captionof{figure}{Elongation systems extend fatty-acid chains beyond the products of fatty-acid synthase.}\labfig{appendix-figure-d4-fatty-acid-elongation-extends-beyond-carbons}\end{center}

\subsubsection{Synthesis of Triacylglycerols}

Triacylglycerols are synthesized in the liver and adipose tissue (specifically in the smooth endoplasmic reticulum) after the acquisition of free fatty acid. In both tissues, triglycerides are produced by a pathway which uses phosphatidic acid as the major intermediate. However, the source of G-3-P is different depending on the type of tissue. In the liver, glycerol is used in the synthesis, as glycerol kinase is present to produce phosphorylated glycerol. Meanwhile, in adipose tissue without glycerol kinase, dihydroxyacetone phosphate, an intermediate of the glycolysis pathway, is used instead as the precursor. In both pathways, 2 acyl-CoA are used to convert glycerol-3-phosphate into phosphatidic acid. The end product of this process are triglycerides. These triglycerides are incorporated into secreted VLDL in the liver.

\begin{center}\includegraphics[width=0.82\textwidth,height=0.68\textheight,keepaspectratio]{\reviewnotesimagepath d4-figure-triacylglycerols-synthesized-liver-adipose-tissue.png}\par\captionof{figure}{Triacylglycerol synthesis and storage in liver and adipose tissue.}\labfig{appendix-figure-d4-triacylglycerols-synthesized-liver-adipose-tissue-specifically}\end{center}

\subsubsection{Cholesterol Synthesis}

\begin{center}\includegraphics[width=0.82\textwidth,height=0.68\textheight,keepaspectratio]{\reviewnotesimagepath d4-figure-cholesterol-synthesis.png}\par\captionof{figure}{Major steps in cholesterol biosynthesis from acetyl-CoA.}\labfig{appendix-figure-d4-cholesterol-synthesis}\end{center}

Malonyl-CoA is the key precursor for the synthesis of cholesterol. It is the intermediate produced from acetoacetyl-CoA by hydroxymethylglutaryl-CoA synthase, which produces beta hydroxymethylglutaryl-CoA reductase (which may be regulated by drugs such as lovastatin to control cholesterol levels), which is later reduced by hydroxymethylglutaryl-CoA reductase into mevalonate. Additionally, cholesterol is used to produce bile salts and bile acids. Examples are shown below.

\begin{center}\includegraphics[width=0.82\textwidth,height=0.68\textheight,keepaspectratio]{\reviewnotesimagepath d4-figure-malonyl-co-precursor-synthesis-cholesterol.png}\par\captionof{figure}{Acetyl-CoA-derived intermediates support fatty-acid and sterol synthesis.}\labfig{appendix-figure-d4-malonyl-co-precursor-synthesis-cholesterol-intermediate-acetoacetyl-co}\end{center}

\subsection{Fatty Acid Metabolism Coordination}

ACC (acetyl-CoA carboxylase) is regulated by glucagon and insulin. It controls the production of malonyl-CoA from acetyl-CoA. Apart from serving a role in the synthesis of fatty acids, malonyl-CoA also acts as an inhibitor of CPT1 in carnitine shuttles (which transport long chain fatty acids into the mitochondria). This means that the synthesis of malonyl-CoA prevents a futile cycle of simultaneous fatty acid degradation during synthesis.

\begin{center}\includegraphics[width=0.82\textwidth,height=0.68\textheight,keepaspectratio]{\reviewnotesimagepath d4-figure-acc-acetyl-co-carboxylase-regulated.png}\par\captionof{figure}{Phosphorylation and allosteric signals regulate acetyl-CoA carboxylase.}\labfig{appendix-figure-d4-acc-acetyl-co-carboxylase-regulated-glucagon-insulin}\end{center}

The main differences between fatty-acid degradation and biosynthesis are summarized in Table~\ref{tab:appendix-table-d4-fatty-acid-degradation-biosynthesis}.
\pagelayout{wide}
\begin{center}
\captionof{table}{Comparison of fatty-acid degradation and biosynthesis.}\labtab{appendix-table-d4-fatty-acid-degradation-biosynthesis}
\small
\renewcommand{\arraystretch}{1.15}
\begin{tabularx}{\textwidth}{XX}
\toprule
\textbf{Degradation} & \textbf{Biosynthesis} \\
\midrule
1. Product is acetyl-CoA. & Precursor is acetyl-CoA. \\

2. Malonyl-CoA is not involved; no requirement for biotin. & Malonyl-CoA is source of two-carbon units; biotin required. \\

3. Oxidative process; requires NAD\textsuperscript{+} and FAD and produces ATP. & Reductive process; requires NADPH and ATP. \\

4. Fatty acids form thioesters with CoA--SH. & Fatty acids form thioesters with acyl carrier proteins (ACP--SH). \\

5. Starts at carboxyl end (\(\mathrm{CH_3CO_2^-}\)). & Starts at methyl end (\(\mathrm{CH_3CH_2^-}\)). \\

6. Occurs in the mitochondrial matrix, with no ordered aggregate of enzymes. & Occurs in the cytosol, catalyzed by an ordered multienzyme complex. \\

7. \(\beta\)-Hydroxyacyl intermediates have the \textit{L} configuration. & \(\beta\)-Hydroxyacyl intermediates have the \textit{D} configuration. \\
\bottomrule
\end{tabularx}
\end{center}
\pagelayout{margin}

\subsection{Required Fatty Acids}

The body does not only require saturated fatty acids; it also uses unsaturated and polyunsaturated fatty acids. As such, there is a mechanism to desaturate fatty acids. This system uses molecular oxygen, NADH and cytochrome b5. This process occurs inside of the endoplasmic reticulum.

In addition to the desaturation mechanism, some fatty acids need to be received from food. These are the linoleic and linolenic (omega 3 and 6) fatty acids, which are referred to as the essential fatty acids. These fatty acids are essential for eicosanoid synthesis, which is used in the production of hormones such as prostaglandins, thromboxanes and leukotrienes.

\subsection{Lipids and the Fasting Metabolism}

During low energy stages, glucagon and epinephrine are activated, which causes the activation of hormone sensitive lipase and the release of fat into the circulation. This fat enters into the liver and provides energy for the function of the liver through beta oxidation and the TCA cycle. In addition, glucose is also produced by liver cells and released into circulation.

During the fatty acid metabolism, there is an increased production of acetyl-COA, which results from beta oxidation. This acetyl-CoA may be changed into ketone bodies (by ketogenesis), which may be used as another source of energy for the body during the starvation stage. Acetyl-CoA is not used in the TCA because oxaloacetate acts as both a precursor for gluconeogenesis and is a part of the TCA cycle. Since gluconeogenesis is preferred, the TCA cycle is not fed.

Overall, this leads to a high level of fat derived ATP and NADH in the liver mitochondria, which inhibits isocitrate dehydrogenase, and shifts the oxaloacetate malate equilibrium towards malate and hence gluconeogenesis. This is summarized below.

\begin{center}\includegraphics[width=0.82\textwidth,height=0.68\textheight,keepaspectratio]{\reviewnotesimagepath d4-figure-overall-leads-high-level-fat.png}\par\captionof{figure}{During fasting, fatty-acid oxidation supports hepatic ketone-body production.}\labfig{appendix-figure-d4-overall-leads-high-level-fat-derived}\end{center}

In addition to the above mentioned effects, the fat derived ATP and NADH acts to inhibit isocitrate dehydrogenase and malate dehydrogenase. This causes further suppression of the krebs cycle and further production of ketone bodies.

Ketone bodies are synthesized from 2 molecules of acetyl-CoA attached together by thiolase, which also catalyzes the removal of one coenzyme A group. This results in acetoacetyl-CoA. Following this, HMG-CoA synthase produces HMG-CoA by removal of coenzyme a and substituting another acetyl-CoA. HMG-CoA lyase then removes one acetyl-CoA and produces acetoacetate, which may be isomerized by hydroxybutyrate dehydrogenase into beta hydroxybutyrate (the main form of ketone body, along with acetoacetate). Alternatively, the removal of a carbon dioxide molecule yields acetone.

\begin{center}\includegraphics[width=0.82\textwidth,height=0.68\textheight,keepaspectratio]{\reviewnotesimagepath d4-figure-ketone-bodies-synthesized-molecules-acetyl.png}\par\captionof{figure}{Ketogenesis converts acetyl-CoA into acetoacetate and beta-hydroxybutyrate.}\labfig{appendix-figure-d4-ketone-bodies-synthesized-molecules-acetyl-co-attached}\end{center}

\begin{center}\includegraphics[width=0.82\textwidth,height=0.68\textheight,keepaspectratio]{\reviewnotesimagepath d4-figure-ketone-bodies-synthesized-molecules-acetyl-02.png}\par\captionof{figure}{Ketogenesis converts acetyl-CoA into acetoacetate and beta-hydroxybutyrate.}\labfig{appendix-figure-d4-ketone-bodies-synthesized-molecules-acetyl-co-attached-02}\end{center}

\begin{center}\includegraphics[width=0.82\textwidth,height=0.68\textheight,keepaspectratio]{\reviewnotesimagepath d4-figure-ketone-bodies-synthesized-molecules-acetyl-03.png}\par\captionof{figure}{Ketogenesis converts acetyl-CoA into acetoacetate and beta-hydroxybutyrate.}\labfig{appendix-figure-d4-ketone-bodies-synthesized-molecules-acetyl-co-attached-03}\end{center}

\section{Protein Metabolism}

Unlike carbohydrate and fat, excess amino acids are not stored as a dedicated reserve. Their amino groups are transferred and removed during catabolism; nitrogen is transported and excreted largely as urea, while carbon skeletons may be oxidized or used in biosynthesis. During fasting or illness, proteolysis can supply amino acids for hepatic gluconeogenesis; insulin generally promotes protein synthesis and suppresses protein breakdown.

On average, a person contains 20\% protein, of which 2\% is broken down and resynthesized daily. The majority of this protein is synthesized from recycled amino acids originating from protein breakdown (75\%). The rest originates from dietary intake of amino acids. The rate of daily energy production may be increased by some disease and/or starvation.

The nitrogen balance is the amount of amino acid consumed versus the amount excreted. Ideally, this is equal. If it is negative, a loss of body mass results, while positive balances lead to the accretion of mass. This is often found in growing children or pregnancy.

Amino acids are synthesized from metabolic intermediates of the pentose phosphate pathway, glycolysis, and the citric acid cycle. Amino acids differ because of their R groups (functional groups), and are classified by whether humans can synthesize them. Table~\ref{tab:appendix-table-d4-amino-acid-requirements} lists the nonessential, conditionally essential, and essential amino acids shown in the source.

\begin{center}
\captionof{table}{Nonessential and essential amino acids for humans and the albino rat.}\labtab{appendix-table-d4-amino-acid-requirements}
\small
\renewcommand{\arraystretch}{1.15}
\begin{tabularx}{\textwidth}{XXX}
\toprule
\textbf{Nonessential} & \textbf{Conditionally essential\textsuperscript{*}} & \textbf{Essential} \\
\midrule
Alanine & Arginine & Histidine \\

Asparagine & Cysteine & Isoleucine \\

Aspartate & Glutamine & Leucine \\

Glutamate & Glycine & Lysine \\

Serine & Proline & Methionine \\

 & Tyrosine & Phenylalanine \\

 & & Threonine \\

 & & Tryptophan \\

 & & Valine \\
\bottomrule
\end{tabularx}
\end{center}
\noindent\textsuperscript{*}Required to some degree in young, growing animals and/or sometimes during illness.\\
\noindent\textit{Source: Lehninger Principles of Biochemistry, 6th ed., Table 18--1, © 2013 W. H. Freeman and Company.}



\textbf{Essential amino acids are the amino acids which cannot be synthesized in humans. They must be consumed in the diet.}

The degradation of amino acids result in two main parts; the carbon skeleton and nitrate group. The carbon skeleton enters the citric acid cycle to be used as energy, while the nitrate group enters the urea cycle to be excreted or reused into new protein.

\begin{center}\includegraphics[width=0.82\textwidth,height=0.68\textheight,keepaspectratio]{\reviewnotesimagepath d4-figure-degradation-amino-acids-parts-carbon.png}\par\captionof{figure}{Amino-acid catabolism channels carbon skeletons into central metabolic pathways.}\labfig{appendix-figure-d4-degradation-amino-acids-parts-carbon-skeleton}\end{center}

\begin{center}\includegraphics[width=0.82\textwidth,height=0.68\textheight,keepaspectratio]{\reviewnotesimagepath d4-figure-degradation-amino-acids-parts-carbon-02.png}\par\captionof{figure}{Amino-acid catabolism channels carbon skeletons into central metabolic pathways.}\labfig{appendix-figure-d4-degradation-amino-acids-parts-carbon-skeleton-02}\end{center}

\subsection{Transdeamination}

The removal of proteins form amino acids occurs in two steps:\\

\begin{center}\includegraphics[width=0.82\textwidth,height=0.68\textheight,keepaspectratio]{\reviewnotesimagepath d4-figure-removal-proteins-amino-acids-steps.png}\par\captionof{figure}{Overview of amino-group removal and carbon-skeleton metabolism.}\labfig{appendix-figure-d4-removal-proteins-amino-acids-steps}\end{center}

\begin{center}\includegraphics[width=0.82\textwidth,height=0.68\textheight,keepaspectratio]{\reviewnotesimagepath d4-figure-removal-proteins-amino-acids-steps-02.png}\par\captionof{figure}{Overview of amino-group removal and carbon-skeleton metabolism.}\labfig{appendix-figure-d4-removal-proteins-amino-acids-steps-02}\end{center}

\begin{center}\includegraphics[width=0.82\textwidth,height=0.68\textheight,keepaspectratio]{\reviewnotesimagepath d4-figure-removal-proteins-amino-acids-steps-03.png}\par\captionof{figure}{Overview of amino-group removal and carbon-skeleton metabolism.}\labfig{appendix-figure-d4-removal-proteins-amino-acids-steps-03}\end{center}

Transamination uses transaminase and is located in the cytoplasm, while oxidative deamination occurs in the mitochondria and is controlled by glutamate dehydrogenase. Both processes occur mainly inside of hepatic tissues. The removal of an amino group converts amino acids to an alpha keto acid, which is equivalent to the carbon skeleton of the amino acid.

\begin{center}\includegraphics[width=0.82\textwidth,height=0.68\textheight,keepaspectratio]{\reviewnotesimagepath d4-figure-transamination-uses-transaminase-located-cytoplasm.png}\par\captionof{figure}{Transamination transfers amino groups between amino acids and alpha-keto acids; glutamate can undergo oxidative deamination.}\labfig{appendix-figure-d4-transamination-uses-transaminase-located-cytoplasm-oxidative}\end{center}

\textbf{Transaminase} (also known as aminotransferase) enzymes are enzymes responsible for the transfer of an amino group from an amino acid to an alpha keto acid. \textbf{All aminotransferases use pyridoxal phosphate (a B group vitamin) as a cofactor.} Upon transfer of an amino group, the amino acid becomes an alpha keto acid, and vice versa. The most common amino acid acceptors are alpha ketoglutarate (forming glutamate) and oxaloacetate (forming aspartate).

The quantity of and ratio between aspartate transaminase (ast) and alanine transaminase (alt) can be used to indicate various diseases.

Viral hepatitis: AST and ALT increase 10-100 folds, some cases with acute hepatitis (ALT>AST)

Cirrhosis: usually chronic AST>ALT (\textasciitilde{}2.5 folds)

Cardiac injury is evaluated using cardiac troponin I or T in the appropriate clinical context. AST is present in many tissues and is not a specific diagnostic marker for myocardial infarction.

Most transaminases have a preference regarding their amino group acceptor, which is either alpha ketoglutarate or oxaloacetate. The majority of amino groups from amino acids end up in glutamate or aspartate. As such, alanine, glutamate and aspartate are the most abundant amino acids, and are central to the amino acid metabolism.

Alpha keto acids are also metabolic intermediates interconnected with the generation of cellular ATP. Oxaloacetate and alpha ketoglutarate are intermediates of the TCA cycle, while oxaloacetate and pyruvate are glucogenic precursors.

\textbf{Oxidative deamination} removes an amino group. Glutamate dehydrogenase catalyzes the reversible conversion of glutamate to alpha-ketoglutarate and ammonium in mitochondria. Its activity responds to cellular energy signals; ammonia is incorporated into urea in the liver or transported safely in forms such as glutamine.

Glucogenic amino acids yield carbon skeletons that can enter gluconeogenesis. Ketogenic amino acids yield acetyl-CoA or acetoacetate, which can support ketone-body or fatty-acid synthesis. Most amino acids are both glucogenic and ketogenic; leucine and lysine are exclusively ketogenic.

\begin{center}\includegraphics[width=0.82\textwidth,height=0.68\textheight,keepaspectratio]{\reviewnotesimagepath d4-figure-glucogenic-amino-acids-skeletons-converted.png}\par\captionof{figure}{Glucogenic amino-acid skeletons enter gluconeogenesis; ketogenic skeletons form acetyl-CoA or ketone bodies.}\labfig{appendix-figure-d4-glucogenic-amino-acids-skeletons-converted-intermediates}\end{center}

\begin{center}\includegraphics[width=0.82\textwidth,height=0.68\textheight,keepaspectratio]{\reviewnotesimagepath d4-figure-glucogenic-amino-acids-skeletons-converted-02.png}\par\captionof{figure}{Glucogenic amino-acid skeletons enter gluconeogenesis; ketogenic skeletons form acetyl-CoA or ketone bodies.}\labfig{appendix-figure-d4-glucogenic-amino-acids-skeletons-converted-intermediates-02}\end{center}

The fate of the ammonium ion is being recycled and excreted as urea due to its toxicity (either directly, as uric acid or urea) through the urea cycle in the liver (which \textbf{receives and breaks down glutamate/glutamine}, as extrahepatic tissues are not able to generate urea), and filtered out by the kidney. Tissues other than the liver use glutamate dehydrogenase or glutamine synthase in order to create the two intermediates sent to the liver.

Alternatively, ammonia may also be reused in the synthesis of nitrogenous compounds such as amino acids or purines.

The glucose-alanine cycle is responsible for the conversion of pyruvate into alanine by transamination of an amino acid. It is a means of disposing of ammonium groups, and passes alanine to the liver for breakdown. Once broken down, urea is generated along with pyruvate, which may be used in gluconeogenesis to generate glucose for use in body tissues.

\begin{center}\includegraphics[width=0.82\textwidth,height=0.68\textheight,keepaspectratio]{\reviewnotesimagepath d4-figure-glucose-alanine-cycle-responsible-conversion.png}\par\captionof{figure}{The glucose–alanine cycle carries nitrogen from muscle to liver and returns carbon for gluconeogenesis.}\labfig{appendix-figure-d4-glucose-alanine-cycle-responsible-conversion-pyruvate-alanine}\end{center}

\subsection{The Urea Cycle}

The urea cycle occurs partly in the mitochondrial matrix and partly in the cytosol of hepatocytes. It converts nitrogen from free ammonia and aspartate into urea. In the mitochondrial matrix, carbamoyl phosphate synthetase I (CPS1) uses bicarbonate, free ammonia, and two ATP to form carbamoyl phosphate; CPS1 is activated by N-acetylglutamate. Ornithine transcarbamylase then combines carbamoyl phosphate with ornithine to form citrulline, which moves into the cytosol. The cycle consumes three ATP molecules (four high-energy phosphate bonds) per urea molecule.

Aspartate supplies the second nitrogen atom in urea and is generated by transamination of oxaloacetate. Ammonia used by CPS1 is released in the mitochondrial matrix, including through amino-acid catabolism. Citrulline leaves the mitochondrion and combines with aspartate in the cytosol to continue the cycle.

\begin{center}\includegraphics[width=0.82\textwidth,height=0.68\textheight,keepaspectratio]{\reviewnotesimagepath d4-figure-mitochondria-responsible-converting-glutamine-glutamate.png}\par\captionof{figure}{Mitochondrial nitrogen metabolism supplies ammonia and aspartate to the hepatic urea cycle.}\labfig{appendix-figure-d4-mitochondria-responsible-converting-glutamine-glutamate-aspartate}\end{center}

\begin{center}\includegraphics[width=0.82\textwidth,height=0.68\textheight,keepaspectratio]{\reviewnotesimagepath d4-figure-mitochondria-responsible-converting-glutamine-glutamate-02.png}\par\captionof{figure}{Mitochondrial nitrogen metabolism supplies ammonia and aspartate to the hepatic urea cycle.}\labfig{appendix-figure-d4-mitochondria-responsible-converting-glutamine-glutamate-aspartate-02}\end{center}

\begin{center}\includegraphics[width=0.82\textwidth,height=0.68\textheight,keepaspectratio]{\reviewnotesimagepath d4-figure-mitochondria-responsible-converting-glutamine-glutamate-03.png}\par\captionof{figure}{Mitochondrial nitrogen metabolism supplies ammonia and aspartate to the hepatic urea cycle.}\labfig{appendix-figure-d4-mitochondria-responsible-converting-glutamine-glutamate-aspartate-03}\end{center}

Carbamoyl phosphate synthase is allosterically activated by N-acetylglutamate, which is a product synthesized from acetyl-CoA and glutamate when concentrations of both are high (N-acetylglutamate synthase is activated by arginine). In all, high levels of amino acids upregulate the urea cycle.

\subsection{Urea Cycle Disorders}

Urea-cycle disorders impair nitrogen disposal and can cause hyperammonemia and neurological toxicity. Severity varies by disorder and may range from neonatal crises to later-onset episodes.

\begin{center}\includegraphics[width=0.82\textwidth,height=0.68\textheight,keepaspectratio]{\reviewnotesimagepath d4-figure-since-there-no-alternate-pathway.png}\par\captionof{figure}{The hepatic urea cycle converts nitrogen waste into urea for excretion.}\labfig{appendix-figure-d4-since-there-no-alternate-pathway-urea}\end{center}

Hyperammonemia is an elevated blood ammonia concentration; reference ranges vary with age, specimen handling, and laboratory. Severe or rapidly rising concentrations can cause encephalopathy, vomiting, seizures, and cerebral edema. Causes include advanced liver disease and inherited defects in the urea cycle. Ornithine transcarbamylase deficiency is X-linked; males are often severely affected, while heterozygous females can have variable manifestations.

\begin{center}\includegraphics[width=0.82\textwidth,height=0.68\textheight,keepaspectratio]{\reviewnotesimagepath d4-figure-hyperammonemia-blood-ammonia-level-m.png}\par\captionof{figure}{Clinical effects of elevated blood ammonia; diagnostic thresholds vary by age and laboratory.}\labfig{appendix-figure-d4-hyperammonemia-blood-ammonia-level-m-normal}\end{center}

\subsection{Urea Cycle Diseases}

(Credit P’s!)

Example disorders (don’t focus here except albinism, maple syrup urine and phenylketonuria)

The source table includes examples and approximate incidence values for inherited amino-acid catabolism disorders (Table~\ref{tab:appendix-table-d4-amino-acid-catabolism-disorders}).
\pagelayout{wide}
\begin{center}
\captionof{table}{Selected inherited disorders affecting amino-acid catabolism.}\labtab{appendix-table-d4-amino-acid-catabolism-disorders}
\footnotesize
\renewcommand{\arraystretch}{1.15}
\begin{tabularx}{\textwidth}{p{0.19\textwidth}p{0.12\textwidth}p{0.18\textwidth}p{0.21\textwidth}X}
\toprule
\textbf{Medical condition} & \textbf{Approximate incidence (per 100,000 births)} & \textbf{Defective process} & \textbf{Defective enzyme} & \textbf{Symptoms and effects} \\
\midrule
Albinism & $<3$ & Melanin synthesis from tyrosine & Tyrosine 3-monooxygenase (tyrosinase) & Lack of pigmentation; white hair, pink skin \\

Alkaptonuria & $<0.4$ & Tyrosine degradation & Homogentisate 1,2-dioxygenase & Dark pigment in urine; late-developing arthritis \\

Argininemia & $<0.5$ & Urea synthesis & Arginase & Intellectual disability \\

Argininosuccinic acidemia & $<1.5$ & Urea synthesis & Argininosuccinase & Vomiting; convulsions \\

Carbamoyl phosphate synthetase 1 deficiency & $<0.5$ & Urea synthesis & Carbamoyl phosphate synthetase 1 & Lethargy; convulsions; early death \\

Homocystinuria & $<0.5$ & Methionine degradation & Cystathionine \(\beta\)-synthase & Faulty bone development; intellectual disability \\

Maple syrup urine disease (branched-chain ketoaciduria) & $<0.4$ & Isoleucine, leucine, and valine degradation & Branched-chain \(\alpha\)-keto acid dehydrogenase complex & Vomiting; convulsions; intellectual disability; early death \\

Methylmalonic acidemia & $<0.5$ & Conversion of propionyl-CoA to succinyl-CoA & Methylmalonyl-CoA mutase & Vomiting; convulsions; intellectual disability; early death \\

Phenylketonuria & $<8$ & Conversion of phenylalanine to tyrosine & Phenylalanine hydroxylase & Neonatal vomiting; intellectual disability \\
\bottomrule
\end{tabularx}
\par\smallskip
\begin{minipage}{\textwidth}
\footnotesize\raggedright
Source: Lehninger Principles of Biochemistry, 5th ed., Table 18--2, © 2008 W. H. Freeman and Company. The source uses an outdated term for intellectual disability.
\end{minipage}
\end{center}
\pagelayout{margin}

\textbf{Hyperammonemia}

\begin{itemize}

\item A metabolic condition characterized by raised levels of ammonia (>60 μM)

\item Symptoms: neurotoxic CNS, tremor, slurred speech, drowsiness, vomiting, cerebral edema, blurred vision

\item Acquired hyperammonemia: caused by liver cirrhosis and liver diseases (e.g., hepatitis B, hepatotoxin)

\item Congenital hyperammonemia: caused by genetic defects in the enzymes of the urea cycle

\begin{itemize}

\item Ornithine transcarbamylase deficiency is X-linked; males are often more severely affected, while heterozygous females may have variable manifestations.

\end{itemize}

\end{itemize}

Laboratory patterns associated with several urea-cycle enzyme deficiencies are summarized in Table~\ref{tab:appendix-table-d4-urea-cycle-deficiencies}.
\pagelayout{wide}
\begin{center}
\captionof{table}{Selected laboratory findings in urea-cycle enzyme deficiencies.}\labtab{appendix-table-d4-urea-cycle-deficiencies}
\small
\renewcommand{\arraystretch}{1.15}
\begin{tabularx}{\textwidth}{p{0.14\textwidth}p{0.08\textwidth}p{0.24\textwidth}p{0.24\textwidth}X}
\toprule
\textbf{Location} & \textbf{Abb.} & \textbf{Enzyme} & \textbf{Disorder} & \textbf{Measurements} \\
\midrule
Mitochondria & CPS1 & Carbamoyl phosphate synthetase I & Carbamoyl phosphate synthetase I deficiency & +Ammonia \\

Mitochondria & OTC & Ornithine transcarbamylase & Ornithine transcarbamylase deficiency & +Ornithine, +Uracil, +Orotic acid \\

Cytosol & ASS & Argininosuccinic acid synthetase & Citrullinemia & +Citrulline \\

Cytosol & ASL & Argininosuccinic acid lyase & Argininosuccinic acid aciduria & +Citrulline, +Argininosuccinic acid \\

Cytosol & ARG & Arginase & Argininemia & +Arginine \\
\bottomrule
\end{tabularx}
\par\smallskip
\begin{minipage}{\textwidth}
\footnotesize\raggedright
Source note: The source image overlays “Argininosuccinase” in red on the argininosuccinic acid lyase row.
\end{minipage}
\end{center}
\pagelayout{margin}

\textbf{Phenylketonuria (PKU)}

\begin{itemize}

\item Defective process: conversion of phenylalanine (aromatic amino acid) to tyrosine

\item Defective enzyme: phenylalanine hydroxylase (requires tetrahydrobiopterin BH4 as a cofactor)

\begin{itemize}

\item Phenylalanine is converted to phenylpyruvate instead which accumulates in tissues, blood, and urine

\end{itemize}

\item Symptoms may include neonatal vomiting, developmental delay or intellectual disability, progressive neurological injury, seizures, and a musty odor of sweat or urine.

\item Detected by biochemical screening (ferric chloride test which detects phenylpyruvate) in newborn babies \(\rightarrow\) individual will require lifelong dietary phenylalanine restriction and tyrosine supplementation (because phenylalanine cannot be converted to tyrosine)\\

\end{itemize}

\textbf{Maple Syrup Urine Disease}

\begin{itemize}

\item \textbf{Defective Process}: Degradation of branched-chain amino acids (isoleucine, leucine, and valine).

\item \textbf{Defective Enzyme}: Branched-chain α-keto acid dehydrogenase complex.

\begin{itemize}

\item Breakdown of branched-chain amino acids (BCAAs) is initiated in the muscle by BCAA aminotransferase.

\item BCAA aminotransferase catalyzes the conversion of BCAAs and α-ketoglutarate to branched-chain α-keto acid and glutamate (e.g., leucine \(\rightarrow\) α-ketoisocaproate).

\item Degradation of branched-chain α-keto acid intermediate is catalyzed by the branched-chain α-keto acid dehydrogenase (BCKD) complex.

\item Genetic deficiency (autosomal recessive inheritance) of the BCKD complex causes maple syrup urine disease.

\end{itemize}

\item \textbf{Symptoms and effects}: Vomiting, seizures, developmental impairment, and a characteristic sweet odor of urine; untreated severe disease can be fatal.

\end{itemize}

\begin{center}\includegraphics[width=0.82\textwidth,height=0.68\textheight,keepaspectratio]{\reviewnotesimagepath d4-figure-symptoms-effects-vomiting-convulsions-mental.png}\par\captionof{figure}{Neurological and gastrointestinal manifestations of maple syrup urine disease.}\labfig{appendix-figure-d4-symptoms-effects-vomiting-convulsions-mental-retardation}\end{center}

\textbf{Treatment of Urea Cycle Disorders:}

\begin{itemize}

\item Restrict dietary protein intake

\item Supplement diet with arginine

\item Enrich diet with alpha ketoglutarate

\item Liver transplant

\item Supplement diet with compounds that provide alternative means of nitrogen excretion e.g.

\item benzoate or phenylbutyrate

\end{itemize}

\subsection{Amino Acid Metabolism Regulation}

There are 2 key control points for amino acid metabolism regulation: glutamate dehydrogenase in the mitochondrial matrix, and glutamine synthetase in the cytosol (which act as gateways for nitrogen incorporation into amino acids). These points are subject to allosteric and (for glutamine synthetase) covalent modification, and especially important because of the reversibility of glutamate dehydrogenase during amino acid breakdown and synthesis.

Glutamate dehydrogenase catalyzes the reductive amination of alpha ketoglutarate (adding a nitrate amino group) and oxidative deamination of glutamate (removing a nitrate amino group).

\begin{center}\includegraphics[width=0.82\textwidth,height=0.68\textheight,keepaspectratio]{\reviewnotesimagepath d4-figure-glutamate-dehydrogenase-catalyzes-reductive-amination.png}\par\captionof{figure}{Glutamate dehydrogenase links amino-acid nitrogen metabolism with the citric acid cycle.}\labfig{appendix-figure-d4-glutamate-dehydrogenase-catalyzes-reductive-amination-alpha}\end{center}

Glutamine synthase acts as a key controlling enzyme for the flow of nitrogen into biological molecules, as glutamine serves as an amino group donor for synthesis of amino acid-containing molecules. As such, it is subject to feedback inhibition by 8+ regulators. When all are present, the enzyme activity is fully repressed (allosteric and covalent modification).

\begin{center}\includegraphics[width=0.82\textwidth,height=0.68\textheight,keepaspectratio]{\reviewnotesimagepath d4-figure-glutamine-synthase-acts-controlling-enzyme.png}\par\captionof{figure}{Glutamine synthetase incorporates ammonia into glutamine and supports nitrogen transport.}\labfig{appendix-figure-d4-glutamine-synthase-acts-controlling-enzyme-flow}\end{center}

\begin{center}\includegraphics[width=0.82\textwidth,height=0.68\textheight,keepaspectratio]{\reviewnotesimagepath d4-figure-glutamine-synthase-acts-controlling-enzyme-02.png}\par\captionof{figure}{Glutamine synthetase incorporates ammonia into glutamine and supports nitrogen transport.}\labfig{appendix-figure-d4-glutamine-synthase-acts-controlling-enzyme-flow-02}\end{center}

\begin{center}\includegraphics[width=0.82\textwidth,height=0.68\textheight,keepaspectratio]{\reviewnotesimagepath d4-figure-glutamine-synthase-acts-controlling-enzyme-03.png}\par\captionof{figure}{Glutamine synthetase incorporates ammonia into glutamine and supports nitrogen transport.}\labfig{appendix-figure-d4-glutamine-synthase-acts-controlling-enzyme-flow-03}\end{center}

\section{Nucleotide Structure and Metabolism}

\subsection{Nucleotide structure and function}

Nucleotides consist of a nitrogenous base, a pentose sugar and one to three phosphate groups. Nucleosides are the nucleotides not including a phosphate group. The pentose sugar ribose and deoxyribose differ in that one oxygen is not attached in deoxyribose, leading to higher stability because it is not a hydroxyl group.

\begin{center}\includegraphics[width=0.82\textwidth,height=0.68\textheight,keepaspectratio]{\reviewnotesimagepath d4-figure-nucleotides-consist-nitrogenous-base-pentose.png}\par\captionof{figure}{Nucleotide structure includes a nitrogenous base, pentose sugar, and phosphate group.}\labfig{appendix-figure-d4-nucleotides-consist-nitrogenous-base-pentose-sugar}\end{center}

\begin{center}\includegraphics[width=0.82\textwidth,height=0.68\textheight,keepaspectratio]{\reviewnotesimagepath d4-figure-nucleotides-consist-nitrogenous-base-pentose-02.png}\par\captionof{figure}{Nucleotide structure includes a nitrogenous base, pentose sugar, and phosphate group.}\labfig{appendix-figure-d4-nucleotides-consist-nitrogenous-base-pentose-sugar-02}\end{center}

The nitrogenous bases are heterocyclic organic compounds which contain nitrogen. Purines (adenine and guanine) and  have a 9 member ring structure made of 4 nitrogen and 5 carbon. Hypoxanthine and xanthine are intermediates in the formation of adenine and guanine nucleotides, and uric acid is the end product of purine nucleotide catabolism.

Pyrimidines are made of a 6 member ring of 2 nitrogen and 4 carbon. The pyrimidine bases present in nucleotides are cytosine, uracil and thymine. Orotic acid is intermediates in the formation of pyrimidine nucleotides.

Nucleosides are named as derivatives of bases (by adding -osine instead). An exception is hypoxanthine, which is the inosine nucleoside, and thiamine which has the base deoxythymidine.\\

\subsection{Biosynthesis of purine and pyrimidine bases}

\begin{center}\includegraphics[width=0.82\textwidth,height=0.68\textheight,keepaspectratio]{\reviewnotesimagepath d4-figure-dietary-nucleic-acid-digestion.png}\par\captionof{figure}{Digestion of dietary nucleic acids into mononucleotides, nucleosides, and bases, followed by absorption and metabolism in intestinal mucosal cells.}\labfig{appendix-figure-d4-dietary-nucleic-acid-digestion}\end{center}

Dietary uptake of purine and pyrimidine bases is minimal, and only approximately 5\% of ingested nucleotides will make it into the circulation, either as the free base or as a nucleoside. De novo synthesis of purines and pyrimidines is thus required.

\begin{center}
\captionof{table}{Comparison of purine and pyrimidine synthesis, degradation, and salvage.}\labtab{appendix-table-d4-purine-pyrimidine-pathways}
\small
\renewcommand{\arraystretch}{1.15}
\begin{tabularx}{\textwidth}{p{0.20\textwidth}XX}
\toprule
\textbf{Pathway} & \textbf{Purines} & \textbf{Pyrimidines} \\
\midrule
De novo synthesis & Nitrogen atoms are derived from amino acids; carbon atoms are derived from carbon dioxide and formate. Six ATP are used per AMP or GMP formed. & Aspartate provides three carbon atoms and one nitrogen atom; the other nitrogen and carbon atoms are derived from glutamine and carbon dioxide. \\

Deoxynucleotide formation & Ribose is reduced to deoxyribose. & A unique pathway forms TTP. \\

Degradation & AMP and GMP are converted to uric acid and excreted. & UMP and CMP are converted to beta-alanine (or beta-aminoisobutyrate from dTMP), ammonia, and carbon dioxide. \\

Salvage & After nucleotide degradation, free purines are recycled into nucleotides. & Nucleotides are degraded and pyrimidines are recycled. \\
\bottomrule
\end{tabularx}
\end{center}

\begin{center}\includegraphics[width=0.82\textwidth,height=0.68\textheight,keepaspectratio]{\reviewnotesimagepath d4-figure-biosynthesis-purine-pyrimidine-bases.png}\par\captionof{figure}{Overview of de novo purine and pyrimidine nucleotide synthesis.}\labfig{appendix-figure-d4-biosynthesis-purine-pyrimidine-bases}\end{center}



\section{Genome Organization}
% CONTENT AUDIT: The Word source has two TODO placeholders under Chromosome Organization; they contain no material to transcribe.

\subsection{Genetic Material}

\subsubsection{Introduction}

Genetic material is differently stored in prokaryotic and eukaryotes. In prokaryotes, it is stored in a nucleoid region without a nuclear membrane and well defined separation. This leads to simultaneous transcription and translation. Meanwhile, in eukaryotes, genetic material is stored inside of a membrane bound nucleus, which acts as a pre translational quality control mechanism.

The concept of genetic material was first discovered when bacterial plasmids were transferred from one strain of bacteria to another (heat killed bacteria was mixed with another alive strain), and the characteristics of the killed strain were found post infection. As such, this principle was named as the transformation principle (even though plasmids were initially identified). This discovery was later verified by the identification of genetic material under the microscope during transformation.

However, some viruses don’t have genetic material. Viruses with RNA as their genetic material which also include DNA intermediates in their replication cycle are called retroviruses (as they use reverse transcriptional mechanisms). Human retroviruses include HIV-1 and HIV-2, the cause of the disease AIDS.

\subsubsection{Nucleic Acids}

There are two main types of nucleic acids; deoxyribonucleic acid and ribonucleic acid. They differ in the presence of a hydroxyl group as opposed to a hydrogen at the 2’ carbon. Nucleic acids are composed of a pentose sugar (5 carbon sugar), phosphate (phosphorus and 4 carbon with a negative charge) and a nitrogenous base. There are two types of bases for nucleic acids: purines (adenine and guanine) and pyrimidines (thiamine, cytosine and uracil). The nitrogenous base is a chemically basic nitrogen containing molecule.

The phosphate group of nucleic acids possess a negative charge, and result in the net negative charge of nucleic acids. The phosphates are joined to the C5 hydroxyl of the ribose or deoxyribose sugar (5’ carbon). Nucleotides without phosphate groups are referred to as nucleosides.

\subsubsection{DNA and RNA}

The polymer of nucleotides is DNA and RNA. DNA consists of two antiparallel (running in opposite directions) strands linked together by intermolecular interactions formed into a double helix.

There are three main forms of DNA. The B form, A form and Z form. Each has different density and functions. The B form has 10 base pairs per turn and is used for cellular DNA. A forms are the condensed form of DNA at 11 base pairs per turn. Z forms are 12 base pairs and twisted left. They are associated with gene regulation.

RNA is a single strand polynucleotide, which is typically shorter than DNA (100s to 1000s of strands long). It replaces thymine with uracil, and uses ribose instead of deoxyribose. RNA, even though single stranded, is able to form short secondary structures with sections linked together, such as hairpins and double helical stem regions.

The state of DNA and RNA at room temperature is as described above. However, it is different when the strands are under stress either from temperature or chemicals such as urea. Once in extreme conditions, it may break down into individual strands and lose helical structure. DNA pairs with itself to form partially denatured strands after the stressor, and but needs to be annealed (repaired) with help from enzyme phosphorylation.

\section{Meiosis and Chromosomal Inheritance}

\subsection{Meiosis}

There are two main types of reproduction for living organisms: asexual and sexual. Asexual reproduction refers to reproduction involving single individuals that pass genes along to their offspring without the need to form gametes and mate with another organism. As such, the offspring of asexual reproduction is genetically identical to the parent.

On the other hand, in sexual reproduction, there are two parents for an individual offspring, meaning that genes are passed down from both parents, and offspring have similar but non identical characteristics to both parents.

The karyotype is an organized display of stained chromosomes, typically arranged in homologous pairs by size, banding pattern, and centromere position. A typical human karyotype has 23 pairs: 22 pairs of autosomes and one pair of sex chromosomes. Each homologous pair usually includes one chromosome inherited from each parent.

Sexual reproduction is the process of fertilization for gamete cells into zygotes (sperm and ova). Since gametes are haploid (1n), the fertilized product is diploid (2n). This zygote then becomes an embryo and subsequently a fetus.

Meiosis is the process by which haploid gamete cells are formed. One cell undergoing meiosis results in 4 haploid daughter cells, each genetically unique. This is because of ‘crossing over.’ The process of meiosis is in essence a repeated version of mitosis where homologous chromosomes are separated from each other. However, there are two unique features. The first is synapsis and crossing over, which occurs during prophase I. It is when homologous chromosomes line up next to each other (synapsis) and form a chiasma (where the two join) to cross over (trade genetic material). The other unique feature of meiosis is reduction division, where the total number of chromosomes is halved, and becomes 1n.

In essence, meiosis 1 separates homologous chromosomes, and meiosis 2 separates sister chromatids.

Three important sources of genetic variation are independent assortment, crossing over, and random fertilization. With independent assortment alone, the number of possible chromosome combinations in a gamete is \(2^n\), where \(n\) is the haploid chromosome number.

Independent chromosomal assortment

\begin{center}\includegraphics[width=0.82\textwidth,height=0.68\textheight,keepaspectratio]{\reviewnotesimagepath d4-figure-independent-chromosomal-assortment.png}\par\captionof{figure}{Independent assortment of homologous chromosome pairs during meiosis.}\labfig{appendix-figure-d4-independent-chromosomal-assortment}\end{center}

\begin{center}\includegraphics[width=0.82\textwidth,height=0.68\textheight,keepaspectratio]{\reviewnotesimagepath d4-figure-independent-chromosomal-assortment-02.png}\par\captionof{figure}{Independent assortment of homologous chromosome pairs during meiosis.}\labfig{appendix-figure-d4-independent-chromosomal-assortment-02}\end{center}

\subsection{Inheritance}

Inheritance is the transmission of genetic information between generations. A gene is a DNA sequence that contributes to a functional product; alleles are alternative versions of a gene or genomic region. Phenotype reflects genetic, environmental, and developmental influences. Many humans carry two alleles at an autosomal locus, but ploidy and inheritance patterns vary by locus and chromosome.

Gregor Mendel studied inheritance in pea plants using self- and cross-fertilization and comparing traits across generations. P denotes the parental generation; F\textsubscript{1}, F\textsubscript{2}, and later subscripts denote successive filial generations.

Mendel's pea-plant experiments established patterns of inheritance for traits with discrete alleles. In simple complete-dominance examples, a dominant allele determines the heterozygote phenotype, while a recessive phenotype requires two recessive alleles. Gametes carry one allele at each autosomal locus.

\begin{center}\includegraphics[width=0.82\textwidth,height=0.68\textheight,keepaspectratio]{\reviewnotesimagepath d4-figure-mendel-discovered-there-certain-heritable.png}\par\captionof{figure}{Mendelian inheritance illustrates segregation of alternative alleles across generations.}\labfig{appendix-figure-d4-mendel-discovered-there-certain-heritable-character}\end{center}

A punnett square is used to cross two plants and identify the phenotype of the offspring. It is a 2 by 2 table where each parent trait is placed in one row or column, and repeated for all traits until each column has two alleles. Standard probability (out of 4 applies).

\subsection{Mendel’s Laws}

The law of dominance and uniformity - Some alleles are dominant while others are recessive; an organism with at least one dominant allele will display the effect of the dominant allele.

The law of segregation - During gamete formation, the two alleles separate randomly, a parent passes only one allele for each trait to each offspring

The law of independent assortment - Each allele in a heterozygous individual has an equal probability of being present in the gamete, as the genes for different traits are inherited independently of each other.

A pedigree diagram is often used to represent inheritance.

\begin{center}\includegraphics[width=0.82\textwidth,height=0.68\textheight,keepaspectratio]{\reviewnotesimagepath d4-figure-pedigree-represent-inheritance.png}\par\captionof{figure}{Pedigree conventions for documenting inheritance patterns in families.}\labfig{appendix-figure-d4-pedigree-often-represent-inheritance}\end{center}

There are some types of inheritance possible given Mendel's theory. See slides for more information.

\begin{enumerate}

\item Autosomal Dominant

\item Autosomal Recessive

\item X-linked Dominant

\item X-linked Recessive

\end{enumerate}

There are also other inheritance patterns that are not considered mendelian. They are:

\begin{enumerate}

\item Codominance

\item Incomplete dominance

\end{enumerate}

Multifactorial diseases are illnesses which are not caused directly by one factor, and require both genetic and environmental influences. For example: polygenic heredity, heart disease, diabetes, cancer, alcoholism, certain mental illnesses such as schizophrenia and bipolar disorder.

\section{Human Genetics}

\begin{center}\includegraphics[width=0.82\textwidth,height=0.68\textheight,keepaspectratio]{\reviewnotesimagepath d4-figure-human-genetics-concepts.png}\par\captionof{figure}{Key terms shown: gene, genome, chromosome, chromatin, and nucleosome}\labfig{appendix-figure-d4-genetics-terminology}\end{center}

Most nuclear genetic material is inherited from both parents, with one chromosome of each homologous pair usually inherited from each. Mitochondrial DNA is typically inherited maternally.

Rosalind Franklin's X-ray diffraction work provided important evidence about DNA's structure; DNA itself was known before those studies. DNA is a sequence of four nucleotide types. Some DNA sequences are transcribed into RNA, and protein-coding genes are translated into proteins; many RNAs are not translated. Crystallin genes are strongly expressed in the lens, illustrating tissue-specific gene expression.

The central dogma summarizes the usual flow of sequence information from DNA to RNA to protein. In eukaryotic cells, protein-coding genes are transcribed in the nucleus; mature messenger RNA is exported to the cytoplasm and translated by ribosomes. Not every RNA transcript encodes a protein.

The regulation of the central dogma is achieved in many ways. One of them is the use of histone modification prior to translation. DNA binds (or unbinds from) to histones, which form a structure called a nucleosome when bound with DNA. This structure is attached to a central chromatin fiber, and may be relaxed or condensed during mitosis. In addition, it also acts as a way to regulate the accessibility of the nuclear DNA. This is because when genetic material is condensed, transcription becomes difficult, if not impossible, to achieve, leading to an absence of mRNA for that specific sequence of nucleotides.

Karyotyping can identify large chromosomal-number and structural changes in dividing somatic cells. A typical human karyotype contains 22 pairs of autosomes and one pair of sex chromosomes; common karyotypes are 46,XX and 46,XY. Karyotyping can detect trisomy 21 (Down syndrome), among other chromosomal conditions.

Other methods of gene expression modification are transcription factors which bind to upstream DNA in order to either promote or inhibit the binding of RNA polymerase, post transcriptional modifications (such as the cap and poly A tail), splicing and nuclear export, and finally translational control of the messenger RNA (such as speed of mRNA degradation and availability/quantity in the nucleus).

Epigenetic mechanisms, including DNA methylation and histone modification, can influence chromatin accessibility and gene expression without changing the DNA sequence. These mechanisms can respond to developmental and environmental signals; specific associations with disease do not by themselves establish causation.

\begin{center}\includegraphics[width=0.82\textwidth,height=0.68\textheight,keepaspectratio]{\reviewnotesimagepath d4-figure-method-modify-expression-genes-epigenetic.png}\par\captionof{figure}{Epigenetic regulation of gene expression through chromatin modification.}\labfig{appendix-figure-d4-method-modify-expression-genes-epigenetic-mechanism}\end{center}

In eukaryotes, transcription occurs in the nucleus and translation in the cytoplasm. Prokaryotes lack a nucleus, so transcription and translation can be coupled in the cytoplasm.

During transcription, RNA polymerase reads the DNA template strand 3′ to 5′ and synthesizes RNA 5′ to 3′. A single-nucleotide variant may or may not alter phenotype; the variant in the beta-globin gene that causes sickle cell disease is one example of a disease-associated substitution. DNA sequencing can identify such variants.

Exons are regions retained in mature RNA after splicing; introns are removed from many eukaryotic transcripts. Noncoding DNA also includes regulatory and intergenic sequences, and not all regulatory elements lie upstream of a gene. The proportions of coding and noncoding sequence depend on how they are defined.

DNA methylation can alter gene expression, often in association with reduced transcription at particular regulatory regions; its effect depends on genomic context and is not a universal on/off switch.

Pharmacogenics refer to the genes involving the effects of drugs.
\end{document}